%%
%% This is LaTeX2e input.
%%

%% The following tells LaTeX that we are using the 
%% style file amsart.cls (That is the AMS article style
%%
\documentclass{amsart}
\usepackage{fancyvrb}
\usepackage{amsmath}
\usepackage{amsfonts}
\usepackage{amssymb}
\usepackage{amsthm}
\usepackage{cases}
\usepackage{caption}
\usepackage{xcolor}
\usepackage{tikz}
\usepackage{centernot}
\usepackage{mathtools}
\usepackage[margin=1in]{geometry}
\usepackage[makeroom]{cancel}
\usepackage{chemarrow}
\usetikzlibrary{decorations.pathreplacing, decorations.pathmorphing, math}
\newcommand{\tikznode}[3][inner sep=0pt]{\tikz[remember
picture,baseline=(#2.base)]{\node(#2)[#1]{$#3$};}}



%\usepackage{ytableau}
%\usepackage{witharrows}

\newtheorem{lemma}{Lemma}[section]
\newtheorem{corollary}[lemma]{Corollary}
\newtheorem{conjecture}[lemma]{Conjecture} 
\newtheorem{proposition}[lemma]{Proposition} 
\newtheorem{theorem}[lemma]{Theorem}
\newtheorem{question}{Question}
%%
%% If some other type is need, say conjectures, then it is constructed
%% by editing and uncommenting the following.
%%

%\newtheorem{conj}[thm]{Conjecture} 


%%% 
%%% The following gives definition type environments (which only differ
%%% from theorem type invironmants in the choices of fonts).  The
%%% numbering is still tied to the theorem counter.
%%% 

\theoremstyle{definition}
\newtheorem{definition}[lemma]{Definition}
\newtheorem{example}[lemma]{Example}
\newtheorem{method}[lemma]{Method}
\allowdisplaybreaks
%%
%% Again more of these can be added by uncommenting and editing the
%% following. 
%%

%\newtheorem{note}[thm]{Note}


%%% 
%%% The following gives remark type environments (which only differ
%%% from theorem type invironmants in the choices of fonts).  The
%%% numbering is still tied to the theorem counter.
%%% 


\theoremstyle{remark}

\newtheorem{remark}[lemma]{Remark}


%%%
%%% The following, if uncommented, numbers equations within sections.
%%% 

\numberwithin{equation}{section}


%%%
%%% The following show how to make definition (also called macros or
%%% abbreviations).  For example to use get a bold face R for use to
%%% name the real numbers the command is \mathbf{R}.  To save typing we
%%% can abbreviate as

\newcommand{\R}{\mathbf{R}}  % The real numbers.

%%
%% The comment after the defintion is not required, but if you are
%% working with someone they will likely thank you for explaining your
%% definition.  
%%
%% Now add you own definitions:
%%

%%%
%%% Mathematical operators (things like sin and cos which are used as
%%% functions and have slightly different spacing when typeset than
%%% variables are defined as follows:
%%%

\DeclareMathOperator{\dist}{dist} % The distance.



%%
%% This is the end of the preamble.
%% 
\title{Triple quantum Schubert calculus}
\author{Matthew J. Samuel}
%\ead{matthematics@gmail.com}

%\affiliation[1]{addressline={31 Cherry Blossom Drive},
%city={Monroe Township},
%postcode={NJ 08831},
%country={USA}}


\DeclareMathOperator{\bpath}{BPath}
\DeclareMathOperator{\Sym}{Sym}
\DeclareMathOperator{\Schub}{Schub}
\DeclareMathOperator{\ESchub}{ESchub}
\DeclareMathOperator{\sch}{\mathfrak{S}}
\DeclareMathOperator{\qsch}{\widetilde{\mathfrak{S}}}


\newcommand{\xrightsquigarrow}[2][]{%
	\mathrel{\mathop{\rightsquigarrow}\limits^{\scriptstyle #2}%
	\ifx\hfuzz#1\hfuzz\else_{\scriptstyle #1}\fi}%
}

\newcommand{\tom}[1]{\ensuremath{\xrightsquigarrow{#1}}}
\newcommand{\fx}{\mathrm{fix}}
%\newcommand{\tomm}[2]{\xrightarrow{#1}_{#2}}
%\counterwithout{equation}{section}
%\setcounter{section}{-1}

\begin{document}

\begin{abstract}
\end{abstract}
\maketitle

{\bf Keywords}: Schubert polynomial, structure constant, Littlewood-Richardson rule, coproduct, double Schubert polynomial

{\bf Mathematics classification codes}: 05E05, 05A05, 05A15, 14N10

{\bf Email}: matthematics@gmail.com

\newcommand{\code}{\mathfrak{c}}
\newcommand{\ducode}{\overline{\mathfrak{c}}}
\newcommand{\wte}{\mathbf{E}}
\newcommand{\qelem}{\widetilde{E}}
\newcommand{\pcomp}[1]{\widehat{#1}}
%$$w_0(n)^{(i)}=w_0(n)w_0(n+1-i)w_0(n-i)$$

%\begin{theorem}
%Let $u\in S_n$ and $1\leq i\leq n$. Then we have
%$$\sch_u(y)=\sum_{u'w_0(n)^{(i)}\tom{n+1-i} uw_0(n)}y_i^{\ell(u',u)}\sch_{u'}(y^{(i)})$$
%where $u'\in S_n$ satisfies $\ell(u'w_0(n)^{(i)})+\ell(u')=\ell(w_0(n)^{(i)})$.
%\end{theorem}
\section{Introduction}



\section{Definitions and statement of results}

\subsection{Quantum double Schubert polynomials and triple quantum Schubert structure constants}
We fix $n\geq 1$ and write
$$x=(x_1,\ldots,x_n),\qquad y=(y_1,\ldots,y_n),\qquad q=(q_1,\ldots,q_{n-1}).$$
For $1\leq a<b\leq n$, set
$$q_{ab}:=q_aq_{a+1}\cdots q_{b-1},$$
and use the cohomological grading $\deg q_i=2$ (equivalently $\deg q_{ab}=2(b-a)$).

Let $s_i$ act on the $y$-variables by swapping $y_i,y_{i+1}$, and define the divided difference operators
$$\partial_i^{(y)}f:=\frac{f-s_i f}{y_i-y_{i+1}}\qquad(1\leq i\leq n-1).$$
For $w\in S_n$, if $w=s_{i_1}\cdots s_{i_r}$ is any reduced word, set
$$\partial_w^{(y)}:=\partial_{i_1}^{(y)}\cdots \partial_{i_r}^{(y)}.$$

Following Kirillov--Maeno \cite{kirillov2000quantum}, define the quantum elementary polynomials $e_j^{(q)}(x_1,\ldots,x_k)$ by
$$\Delta_k(t\mid x_1,\ldots,x_k):=\sum_{j=0}^k t^{k-j}e_j^{(q)}(x_1,\ldots,x_k),$$
where
$$\Delta_k(t\mid x_1,\ldots,x_k)=\det\!\begin{pmatrix}
x_1+t & q_1 & 0 & \cdots & 0\\
-1 & x_2+t & q_2 & \ddots & \vdots\\
0 & -1 & x_3+t & \ddots & 0\\
\vdots & \ddots & \ddots & \ddots & q_{k-1}\\
0 & \cdots & 0 & -1 & x_k+t
\end{pmatrix}.
$$

Let $w_0\in S_n$ be the longest permutation. The top quantum double Schubert polynomial is
$$\qsch_{w_0}(x;y):=\prod_{i=1}^{n-1}\Delta_i(-y_{n-i}\mid x_1,\ldots,x_i),$$
and for $w\in S_n$ we define
$$\qsch_w(x;y):=(-1)^{\ell(w_0) - \ell(w)}\partial_{ww_0}^{(y)}\qsch_{w_0}(x;y).$$
Its specialization
$$\qsch_w(x):=\qsch_w(x;0)$$
is the quantum Schubert polynomial. In particular, at $q=0$ these recover the ordinary double and single Schubert polynomials.

Given their stability property, we may let the number of variables be infinite, indexing the polynomials by $S_\infty$. As a module over $\mathbb{Z}[q][y]$, the quantum double Schubert polynomials form a basis of $\mathbb{Z}[q][x,y]$, and we may write
$$\qsch_u(x;y)\qsch_v(x;z)=\sum_{w,d}q^dc_{u,v}^{w,d}(y;z)\qsch_w(x;y)$$
where $d=(d_1,\ldots,d_{n-1})$ is a weak composition and $q^d=q_1^{d_1}\cdots q_{n-1}^{d_{n-1}}$. The coefficients $c_{u,v}^{w,d}(y;z)$ are the \emph{triple quantum Schubert structure constants}, which are polynomials in $y$ and $z$ with integer coefficients. This product respects the grading where for each $i$, $\deg y_i=\deg z_i=1$ and $\deg q_i=2$. In particular, if $c_{u,v}^{w,d}(y;z)\neq 0$, then $\deg c_{u,v}^{w,d}(y;z)=\ell(u)+\ell(v)-\ell(w)+2\sum_i d_i$. Unlike the usual triple Schubert structure constants, they do not exhibit triangularity in Bruhat order, and indeed we may have $\ell(w) < \ell(u) + \ell(v)$. There is an ordering in which the quantum double Schubert polynomials are triangular, but we will not need this in full generality.

\begin{remark}[$\tom{k}$ and the quantum $k$-Bruhat order] \label{remark:qkbruhat}
The \emph{quantum $k$-Bruhat order} was introduced by Benedetti, Bergeron, Colmenarejo, Saliola, and Sottile \cite{benedetti2025quantum} and studied further by Colmenarejo and Mayers \cite{colmenarejo2026quantum}. That order $\preceq_k^{\mathbf q}$ on $S_\infty[\mathbf q]$ is generated by cover relations
$$u'\;\lessdot_k^{\mathbf q}\;\mathbf q^{\gamma}u't_{ab},\qquad a\leq k<b,$$
each of exactly two kinds: a \emph{classical} cover, with $\gamma=0$ and $\ell(u't_{ab})=\ell(u')+1$, or a \emph{quantum} cover, with $\ell(u')+1=\ell(\mathbf q^{\gamma}u't_{ab})=\ell(u't_{ab})+2(b-a)$, that is $\ell(u't_{ab})=\ell(u')-2(b-a)+1$ and $\gamma=e_a+\cdots+e_{b-1}$. 
These are exactly the two alternatives of condition (4) of Definition \ref{definition:pieri}. Consequently every sequence $t_{a_1b_1},\ldots,t_{a_pb_p}$ witnessing $u\tom{k}w$ is a saturated chain from $u$ to $\mathbf q^{\alpha}w$ in the quantum $k$-Bruhat order, where $\mathbf q^{\alpha}$ is the monomial recorded by $q_k(u,w)$; in particular $u\tom{k}w$ implies $u\preceq_k^{\mathbf q}\mathbf q^{\alpha}w$.

The relation $\tom{k}$ is \emph{not} the full order $\preceq_k^{\mathbf q}$, however. Definition \ref{definition:pieri} imposes the additional requirements that the lower indices $a_1,\ldots,a_p$ be pairwise distinct, so that $\tom{k}$ admits only those saturated chains of a special form. Thus $\tom{k}$ is strictly larger than the covering relation $\lessdot_k^{\mathbf q}$ but strictly smaller than $\preceq_k^{\mathbf q}$: a general chain in the quantum $k$-Bruhat order may reuse a lower index $a_i$, and such chains do not contribute to the Pieri expansion. This is the quantum counterpart of the classical situation, where the Pieri relation of \cite{sottile,samuelmolev} is likewise the restriction of the $k$-Bruhat order to chains with distinct lower indices, and it is exactly this restriction that makes the Pieri coefficients well defined.

What we use below is only that each individual step of a $\tom{k}$ chain is a cover, and hence is governed by the following length-free description of the covers \cite[Proposition~3.1]{benedetti2025quantum} (see also \cite[Proposition~4.2]{colmenarejo2026quantum}): with $i\leq k<j$, the step $u'\lessdot_k^{\mathbf q}\mathbf q^{\gamma}u't_{ij}$ is a classical cover exactly when $u'(i)<u'(j)$ and no $l$ with $i<l<j$ satisfies $u'(i)<u'(l)<u'(j)$, and it is a quantum cover exactly when $u'(i)>u'(j)$ and every $l$ with $i<l<j$ satisfies $u'(j)<u'(l)<u'(i)$. This characterization is an immediate consequence of the standard formula for the length of $u't_{ij}$ (Lemma \ref{lemma:lengthchange} below), and it is the form in which the covers are used throughout Section \ref{section:proofpieri}.
\end{remark}


For $d=(0,\ldots,0)$, this is the ordinary triple Schubert calculus structure constant, and it was proved in \cite{gao2025grahampositivitytripleschubert} that these have the positivity property that they are polynomials in the differences $y_i-z_j$ with nonnegative integer coefficients. We will refer to this as \emph{triple Schubert positivity}. We have the following conjecture:
\begin{theorem}[Triple quantum Schubert positivity]
For all $u,v,w$ and $d$, the structure constants $c_{u,v}^{w,d}(y;z)$ are polynomials in the differences $y_i-z_j$ with nonnegative integer coefficients.
\end{theorem}
\begin{proof}
The proof below is adapted from the manuscript of Tongren Xiao and Yihua Yue \cite{xiaoyuequantumpositivity}, with their permission. Their product is written with the output basis carrying the parameters of the second factor. Commutativity of multiplication therefore identifies their coefficient $C_{v,u}^{w,d}(z,y)$ with $c_{u,v}^{w,d}(y;z)$. It is enough to prove that $C_{u,v}^{w,d}(y,t)$ is in $\mathbb{Z}_{\geq 0}[t_i-y_j]$.

Fix $u,v,w\in S_n$, increasing $n$ if necessary so that all permutations and parameters occurring in the coefficient under consideration are included. Write
$\mathbf a=(t_1,\ldots,t_n,y_1,\ldots,y_n)$, and let $\tau\in S_{2n}$ interchange the two blocks, so that $\tau\mathbf a=(y_1,\ldots,y_n,t_1,\ldots,t_n)$. Regard the product expansion in $S_{2n}$:
\begin{equation}\label{eq:xiaoyue-product}
\qsch_u(x;\tau\mathbf a)\qsch_v(x;\mathbf a)
=\sum_{w\in S_{2n}}C_{u,v}^{w,\tau}(\mathbf a;q)\qsch_w(x;\mathbf a).
\end{equation}
We recall the geometric interpretation of these coefficients. Let $e_i(x)$ and $E_i(x,q)$ denote, respectively, the elementary symmetric polynomials and their quantum analogues in $2n$ variables, and set
\[
J_{\mathbf a}^{2n}=\langle e_i(x)-e_i(\mathbf a):1\leq i\leq 2n\rangle,\qquad
J_{q\mathbf a}^{2n}=\langle E_i(x,q)-e_i(\mathbf a):1\leq i\leq 2n\rangle.
\]
The Fomin--Gelfand--Postnikov deformation map \cite{fomin1997quantum}, extended $\mathbb Z[\mathbf a]$-linearly, induces
\[
\bar\theta:\mathbb Z[x,\mathbf a]/J_{\mathbf a}^{2n}
\longrightarrow \mathbb Z[x,\mathbf a,q]/J_{q\mathbf a}^{2n}.
\]
The classical Borel presentation identifies the source with $H_T^*(X)$, and Kim's presentation \cite{kim1996equivariant} identifies the target with $QH_T^*(X)$, where $X=\mathrm{Fl}_{2n}$. Under the Lam--Shimozono identification \cite{lam2014quantum} of quantum double Schubert polynomials with quantum Schubert classes, \eqref{eq:xiaoyue-product} is the quantum equivariant Littlewood--Richardson expansion for the classes obtained by applying $\bar\theta$ to the first factor. Consequently,
\begin{equation}\label{eq:xiaoyue-eqlr}
C_{u,v}^{w,\tau}(\mathbf a;q)
=\sum_{d\geq0}q^d(\pi_{M_d})^T_*
\left(
\operatorname{ev}_1^*[\tau X_u]^T\,
\operatorname{ev}_2^*[X_v]^T\,
\operatorname{ev}_3^*[Y_w]^T
\right).
\end{equation}
Here $M_d=\overline M_{0,3}(X,d)$, $\operatorname{ev}_i:M_d\to X$ are the evaluation maps, $\pi_{M_d}:M_d\to\mathrm{pt}$ is the structure map, and $X_r$ and $Y_r$ denote the Schubert and opposite Schubert varieties indexed by $r$.

For each effective degree $d$, put
\[
F=(\operatorname{ev}_1,\operatorname{ev}_2):M_d\longrightarrow X\times X,\qquad
I_{u,v}^d=F^{-1}(\tau X_u\times X_v),\qquad
\Gamma_{u,v}^d=(\operatorname{ev}_3)_*[I_{u,v}^d].
\]
The incidence scheme $I_{u,v}^d$ is reduced and locally irreducible, and all its irreducible components have codimension $\ell(u)+\ell(v)$. Indeed, the moduli space $M_d$ is irreducible by Thomsen's theorem \cite{thomsen1998irreducibility}, and Mihalcea's incidence lemma \cite[Lemma~6.3]{mihalcea2006positivity} applies after diagonal translation by $a_0=1\times w_0\in S_{2n}$, where $w_0$ is the longest element in the second block: it identifies $\tau X_u\times X_v$ with $w_0^{(2n)}X_u\times X_v$. The codimension is therefore
\[
\operatorname{codim}_X(w_0^{(2n)}X_u)+\ell(v)=\ell(u)+\ell(v).
\]
By the incidence formula of Mihalcea \cite[Section~3.1]{mihalcea2006positivity}, this gives
\begin{equation}\label{eq:xiaoyue-incidence}
\operatorname{ev}_1^*[\tau X_u]^T\,
\operatorname{ev}_2^*[X_v]^T\cap[M_d]^T=[I_{u,v}^d]_T.
\end{equation}

The cycle $\Gamma_{u,v}^d$ is effective and $B^-(\tau)$-invariant. To see invariance, write $N^-(\tau)=N^-\cap\tau N^-\tau^{-1}$ and $B^-(\tau)=TN^-(\tau)$. The varieties $X_v$ and $\tau X_u$ are stable under the corresponding actions, so their product, its inverse image under the equivariant map $F$, and its proper push-forward under $\operatorname{ev}_3$ are $B^-(\tau)$-invariant. Effectivity follows from the reducedness of $I_{u,v}^d$.

There is a unique expansion
\[
[\Gamma_{u,v}^d]_T=\sum_{r\in S_{2n}}p_{u,v}^{r,d}(\mathbf a)[X_r]_T
\quad\text{in }H_T^*(X).
\]
By Gao and Xiong's Graham--positivity theorem for $B^-(\tau)$-invariant effective cycles \cite[Theorem~2.3 and Corollary~2.4]{gao2025grahampositivitytripleschubert},
\[
p_{u,v}^{r,d}(\mathbf a)\in
\mathbb Z_{\geq0}[t_i-y_j:1\leq i,j\leq n].
\]
Finally, by equivariant Poincaré duality, \eqref{eq:xiaoyue-incidence}, and the projection formula,
\[
(\operatorname{ev}_3)^T_*
\left(\operatorname{ev}_1^*[\tau X_u]^T
\operatorname{ev}_2^*[X_v]^T\right)
=[\Gamma_{u,v}^d]_T,
\]
and hence the degree-$d$ summand of \eqref{eq:xiaoyue-eqlr} is
\[
(\pi_X)^T_*\bigl([\Gamma_{u,v}^d]_T[Y_w]^T\bigr)
=\sum_{r\in S_{2n}}p_{u,v}^{r,d}(\mathbf a)
(\pi_X)^T_*\bigl([X_r]^T[Y_w]^T\bigr)
=p_{u,v}^{w,d}(\mathbf a).
\]
Thus every coefficient $C_{u,v}^{w,d}(y,t)$ is in
$\mathbb Z_{\geq0}[t_i-y_j]$. Replacing $(u,v,y,t)$ by $(v,u,z,y)$ proves the asserted positivity of $c_{u,v}^{w,d}(y;z)$.
\end{proof}

\begin{example} \label{example:cppfpos}
	We compute the product $\qsch_{15423}(x;y)\qsch_{1432}(x;z)$:
\begin{align*}
&q_{2} q_{3} \left(y_{1} + y_{4} - z_{1} - z_{2}\right) \qsch_{124635}{\left(x; y \right)} + q_{2} q_{3} \left(y_{1} + y_{4} - z_{1} - z_{2}\right) \qsch_{13452}{\left(x; y \right)} \\
&+ q_{2} q_{3} \left(y_{1} + y_{4} - z_{1} - z_{2}\right) \qsch_{14253}{\left(x; y \right)} \\
&+ q_{2} q_{3} \qsch_{125634}{\left(x; y \right)} + q_{2} q_{3} \qsch_{126435}{\left(x; y \right)} + q_{2} q_{3} \qsch_{13542}{\left(x; y \right)} + q_{2} q_{3} \qsch_{214635}{\left(x; y \right)} \\
&+ q_{2} q_{3} \qsch_{23451}{\left(x; y \right)} + q_{2} q_{3} \qsch_{24153}{\left(x; y \right)} + q_{2} q_{3} \qsch_{31452}{\left(x; y \right)} + q_{2} q_{3} \qsch_{41253}{\left(x; y \right)} \\
&+ q_{2} \left(y_{1} + y_{4} - z_{1} - z_{2}\right) \qsch_{146235}{\left(x; y \right)} + q_{2} \left(y_{4} + y_{5} - z_{1} - z_{2}\right) \qsch_{24513}{\left(x; y \right)} \\
&+ q_{2} \qsch_{246135}{\left(x; y \right)} + q_{3} \qsch_{1724356}{\left(x; y \right)} + q_{3} \qsch_{261435}{\left(x; y \right)} + q_{3} \qsch_{35142}{\left(x; y \right)} \\
&+ \left(q_{2} q_{3} + q_{2} \left(\left(y_{1} - z_{2}\right) \left(y_{4} + y_{5} - z_{1} - z_{2}\right) + \left(y_{4} - z_{1}\right) \left(y_{5} - z_{1}\right)\right)\right) \qsch_{14523}{\left(x; y \right)} \\
&+ \left(\left(y_{4} - z_{1}\right) \left(y_{5} - z_{3}\right) + \left(y_{4} - z_{1}\right) \left(y_{1} + y_{2} - z_{1} - z_{2}\right) \right. \\
&\quad \left. + \left(y_{5} - z_{2}\right) \left(y_{5} - z_{3}\right) + \left(y_{5} - z_{2}\right) \left(y_{1} + y_{2} - z_{1} - z_{2}\right)\right) \qsch_{25413}{\left(x; y \right)} \\
&+ \left(\left(y_{1} - z_{2}\right) \left(y_{1} - z_{3}\right) \left(y_{4} - z_{2}\right) + \left(y_{1} - z_{2}\right) \left(y_{1} - z_{3}\right) \left(y_{5} - z_{1}\right) + \left(y_{1} - z_{2}\right) \left(y_{5} - z_{1}\right) \left(y_{5} - z_{2}\right) \right. \\
&\quad \left. + \left(y_{1} - z_{3}\right) \left(y_{4} - z_{1}\right) \left(y_{5} - z_{1}\right) + \left(y_{4} - z_{1}\right) \left(y_{5} - z_{1}\right) \left(y_{5} - z_{2}\right)\right) \qsch_{15423}{\left(x; y \right)} \\
&+ \left(q_{2} + \left(y_{1} - z_{2}\right) \left(y_{1} - z_{3}\right) + \left(y_{1} - z_{3}\right) \left(y_{5} - z_{1}\right) + \left(y_{5} - z_{1}\right) \left(y_{5} - z_{2}\right)\right) \qsch_{156234}{\left(x; y \right)} \\
&+ \left(q_{3} + \left(y_{1} - z_{1}\right) \left(y_{4} - z_{1}\right) + \left(y_{5} - z_{2}\right) \left(y_{1} + y_{4} - z_{1} - z_{2}\right)\right) \qsch_{45123}{\left(x; y \right)} \\
&+ \left(q_{2} q_{3} + q_{3} \left(\left(y_{1} - z_{1}\right) \left(y_{1} + y_{5} - z_{2} - z_{3}\right) + \left(y_{5} - z_{2}\right) \left(y_{5} - z_{3}\right)\right)\right) \qsch_{15243}{\left(x; y \right)} \\
&+ q_{3} \left(y_{1} + y_{5} + y_{6} - z_{1} - z_{2} - z_{3}\right) \qsch_{162435}{\left(x; y \right)} + q_{3} \left(y_{1} + y_{2} + y_{5} - z_{1} - z_{2} - z_{3}\right) \qsch_{25143}{\left(x; y \right)} \\
&+ q_{2} q_{3} \left(y_{1} + y_{2} + y_{5} - z_{1} - z_{2} - z_{3}\right) \qsch_{12543}{\left(x; y \right)} + q_{2} q_{3} \left(y_{1} + y_{2} + y_{5} - z_{1} - z_{2} - z_{3}\right) \qsch_{21453}{\left(x; y \right)} \\
&+ q_{2} q_{3} \left((y_{1} - z_{2}) (y_{1} - z_{3}) + (y_{1} - z_{2}) (y_{2} - z_{2}) + (y_{1} - z_{2}) (y_{4} - z_{1}) \right. \\
&\quad \left. + (y_{1} - z_{2}) (y_{5} - z_{1}) + (y_{2} - z_{1}) (y_{4} - z_{1}) + (y_{4} - z_{1}) (y_{5} - z_{3})\right) \qsch_{12453}{\left(x; y \right)} \\
&+ \left(y_{1} + y_{4} - z_{1} - z_{2}\right) \qsch_{1742356}{\left(x; y \right)} + \left(y_{1} + y_{4} - z_{1} - z_{2}\right) \qsch_{461235}{\left(x; y \right)} \\
&+ \left(y_{4} + y_{5} - z_{1} - z_{2}\right) \qsch_{35412}{\left(x; y \right)} + \left(y_{4} + y_{5} - z_{1} - z_{2}\right) \qsch_{45213}{\left(x; y \right)} \\
&+ \left(\left(y_{1} - z_{2}\right) \left(y_{1} - z_{3}\right) + \left(y_{1} - z_{2}\right) \left(y_{5} + y_{6} - z_{1} - z_{2}\right) \right. \\
&\quad \left. + \left(y_{1} - z_{3}\right) \left(y_{4} - z_{1}\right) + \left(y_{4} - z_{1}\right) \left(y_{5} + y_{6} - z_{1} - z_{2}\right)\right) \qsch_{164235}{\left(x; y \right)} \\
&+ \left(y_{1} + y_{2} + y_{5} - z_{1} - z_{2} - z_{3}\right) \qsch_{256134}{\left(x; y \right)} + \left(y_{1} + y_{5} + y_{6} - z_{1} - z_{2} - z_{3}\right) \qsch_{165234}{\left(x; y \right)} \\
&+ \left(y_{1} + y_{2} + y_{4} + y_{5} + y_{6} - 2 z_{1} - 2 z_{2} - z_{3}\right) \qsch_{264135}{\left(x; y \right)} \\
&+ \qsch_{1752346}{\left(x; y \right)} + \qsch_{265134}{\left(x; y \right)} + \qsch_{2741356}{\left(x; y \right)} \\
&+ \qsch_{356124}{\left(x; y \right)} + \qsch_{364125}{\left(x; y \right)} + \qsch_{462135}{\left(x; y \right)} \\
&+ \qsch_{561234}{\left(x; y \right)}
\end{align*}
By inspection, we can see that the coefficients exhibit triple Schubert poisitivity.
\end{example}


We will provide formulas below that demonstrate this in special cases. In particular, we prove

\begin{theorem}
Suppose $u,v,w\in S_\infty$ and $v$ is $k$-Grassmannian for some $k$. Then for any $d$, $c_{u,v}^{w;d}(y;z)$ exhibits triple Schubert positivity.
\end{theorem}

\subsection{The parabolic case}

Unlike ordinary cohomology or equivariant cohomology, quantum cohomology and equivariant quantum cohomology do not demonstrate simple functoriality, in the sense that a projection of the complete flag variety $G/B$ to a partial flag variety $G/P$ does not induce a homomorphism of rings $QH_T^*(G/P)\to QH_T^*(G/B)$, and indeed the quantum cohomology of $G/P$ is not isomorphic to the subring of the quantum cohomology of $G/B$ generated by the Schubert classes indexed by $W^P$. This is reflected in the fact that there are parabolic quantum double Schubert polynomials defined in a different manner that represent Schubert classes in $QH_T^*(G/P)$ \cite{lam2014quantum} and are not in general equal to the quantum double Schubert polynomials above.

\begin{definition}[Parabolic quantum double Schubert polynomials]
	Following Lam and Shimozono \cite{lam2014quantum}, we define parabolic quantum double Schubert polynomials as follows. Fix an integer $n>0$. A parabolic subgroup $W_P$ of $S_{n}$, which has Coxeter generators $s_1,\ldots,s_{n-1}$, is generated by a subset $P\subseteq \{s_1,\ldots,s_{n-1}\}$. In terms of the subgroup's properties, it is most natural to identify $P$ by the reflections that it \emph{excludes}. If $P$ is generated by $S\setminus \{s_{i_1},\ldots,s_{i_k}\}$ in increasing order, we define integers $n_1,\ldots,n_{k+1}$ with
	$$n_1+\cdots+n_{k+1}=n$$
	by $n_1 = i_1$ and $n_p = i_{p} - i_{p-1}$ for $2\leq p\leq k+1$. Then $N_j$ for $j\geq 1$ is defined by
	$$N_j = n_1+\cdots+n_j$$
	for $0\leq j\leq k+1$.
	
	For such a setup, define a matrix $D^P_j$ to be the $n\times n$ matrix with $x_1,\ldots,x_n$ on the diagonal, $-1$ on the superdiagonal, and $(-1)^{n_{j+1}+1}q_j$ at $(N_{j+1},N_{j-1}+1)$ for $1\leq j\leq k$. Notice that there are precisely $k$ $q$ variables, and we set
	$$\deg q_j = n_j + n_{j+1}$$  
	
	Define 
	$$\Delta^P_j(t;x)=\det(D^P_j-t\mathrm{Id})$$
	Then we define the parabolic quantum double Schubert polynomials $\qsch_w^P(x;y)$ for $w\in W^P$ by
	$$\qsch^P_w(x;y)=(-1)^{\ell(w(w_0^P)^{-1})}\partial_{w(w_0^P)^{-1}}^{(y)}\prod_{j=1}^k \Delta^P_j(-y_{n - N_{j+1}+1},\ldots,-y_{n-N_j};x)$$
	where $w_0^P$ is the longest element of $W^P$. 
	
	Taken as defined, when we vary $n$ these polynomials are not stable under inclusion and do not form a basis of $\mathbb{Z}[q][x_1,\ldots,x_n;y]$ as a module over $\mathbb{Z}[q][y]$. To obtain stability and a basis, instead we fix $n$ and choose $N>n$ as large as we would like, maintaining $N_{k+1}=n$ and setting $n_{k+2} = N-n$. From a cohomological perspective, this introduces ``extraneous'' polynomials that are indexed by elements not in $W^P$. The full indexing set is $W^{\pcomp{P}}\subset S_\infty$, where
	$$\pcomp{P} = P \cup \{s_{n+1},s_{n+2},\ldots\}$$
\end{definition}

\begin{example}
As equivariant cohomology rings of finite dimensional projective varieties over $\mathbb{C}$ are finitely generated modules over the coefficient ring, it is fully expected that products of double Schubert polynomials may include terms indexed by permutations that do not lie within the indexing set of the basis of Schubert classes. For example, consider the following non-quantum double Schubert polynomial product, with $P=\{s_1,s_3\}$ and basis elements in $S_4$:
\begin{align*}
\sch_{2413}(x;y)\sch_{3412}(x;y)
&= - (\alpha_1+\alpha_2)(\alpha_1+\alpha_2+\alpha_3)(\alpha_2+\alpha_3)\,\sch_{3412}(x;y) \\
&\quad + (\alpha_1+\alpha_2)(\alpha_1+2\alpha_2+2\alpha_3+\alpha_4)\,\sch_{35124}(x;y) \\
&\quad - (\alpha_1+\alpha_2)\,\sch_{361245}(x;y) \\
&\quad - (\alpha_1+2\alpha_2+2\alpha_3+\alpha_4)\,\sch_{45123}(x;y) \\
&\quad + \sch_{461235}(x;y).
\end{align*}
Not only does the product $\qsch_{2413}^{P}(x;y)\qsch_{3412}^{P}(x;y)$ include extra terms indexed by permutations in $W^P$, there are also terms with nonzero coefficients that are in $W^{\pcomp{P}}\setminus W^P$:
\begin{align*}
q_1(\alpha_1\alpha_2+\alpha_1\alpha_3+\alpha_2\alpha_3+\alpha_2^2)\,\qsch_{132}^{P}(x;y)
&+ q_1(\alpha_1+\alpha_2)\,\qsch_{231}^{P}(x;y) \\
&+ q_1(\alpha_2+\alpha_3)\,\qsch_{1423}^{P}(x;y) \\
&+ \color{red}{q_1(\alpha_1+\alpha_2)\,\qsch_{13254}^{\pcomp{P}}(x;y)} \\
&+ \color{red}{q_1\,\qsch_{14253}^{\pcomp{P}}(x;y)} \\
&+ q_1\,\qsch_{15234}^{P}(x;y) \\
&+ q_1\,\qsch_{2413}^{P}(x;y)
\end{align*}
\end{example}


\begin{definition}[$q$-vectors and Peterson-Woodward lifts]
	Define $Q^\vee$, for legacy purposes, to be the set of weak compositions. For an integer $i$ and $d\in Q^\vee$, define
	$$\langle d\rangle_i = 2d_i - d_{i-1} - d_{i + 1}$$
	The double indexed version will be denoted by
	$$\langle d\rangle_{ij} = \sum_{k=i}^{j-1}\langle d\rangle_k$$

	For a parabolic subgroup $W_P$, we define $Q^\vee_P$ to be the set of weak compositions $d$ such that $\langle d\rangle_i = 0$ for all $s_i\notin P$. The $q$-vectors in $W^P$ are indexed by $Q^\vee/Q^\vee_P$.

	Let $d_P\in Q^\vee/Q^\vee_P$. A \emph{Peterson-Woodward lift} of $d_P$ is a $d\in Q^\vee$ such that $d_P\in d+ Q_P^\vee$ and $\langle d\rangle_{ij}\in \{0,-1\}$ for all $i,j$ such that $s_{i},\ldots,s_{j-1}\in P$.
 \end{definition}

 \begin{method}[Checking the Peterson-Woodward lift property]
 While backing out a Peterson-Woodward lift is difficult, going in the other direction is not. Suppose $d\in Q^\vee$ and $P$ is parabolic. If there is any $i\in P$ such that $\langle d\rangle_i\notin\{0,-1\}$, then $d$ is not a Peterson-Woodward lift. If this condition does hold, let $z_P(d)$ be the set of all indexes $i$ such that $\langle d\rangle_i = 0$. This set is naturally divided into maximal disjoint intervals $I_1,\ldots,I_m$. We must check the additional condition that for each such $I$ and each pair $i<j$ with $i,j-1\in I$, the condition $\langle d\rangle_{ij}\in\{0,-1\}$ holds. If and only if it does, then $d$ is a Peterson-Woodward lift of a unique $d_P\in Q^\vee/Q^\vee_P$. We denote this $d_P$ by $\pi_P(d)$. 
 \end{method}

 This algorithm is polynomial in the length of $d$, which in almost all currently feasible cases is quite small. In contrast, a naive search for $d$ with $d_P$ as the input has an exponential search space, though a more refined algorithm may exist.

 We prove the following result.

 \begin{proposition}
Let $W_P$ be a parabolic subgroup of $S_\infty$ with $P=\pcomp{P}$ and let $u,v,w\in W^{P}$. Then for any $d_P = \pi_P(d)\in Q^\vee/Q_P^\vee$, we have that
$$c_{u,v}^{w,d_P}(y;z) = c_{u,v}^{ww_Pw_{z_P(d)},d}(y;z),$$
\end{proposition}

Thus, though we lack a simple functoriality property, we can reduce the computation of parabolic triple quantum Schubert structure constants completely to the complete flag variety case with an efficient algorithm. We obtain a nontrivial corollary from this.

\begin{example}
	The product in the example above, with $P=\pcomp{P}$ satisfying $P\cap\{s_1,s_2,s_3,s_4,s_5\} = \{s_1,s_4\}$, allows us to compute the triple structure constants for the corresponding parabolic quantum double Schubert polynomials:
\begin{align*}
&q_{1} q_{2} \left(y_{1} + y_{2} + y_{5} - z_{1} - z_{2} - z_{3}\right) \qsch_{1243}^P{\left(x; y \right)} \\
&+ q_{1} q_{2} \qsch_{124365}^P{\left(x; y \right)} + q_{1} q_{2} \qsch_{1342}^P{\left(x; y \right)} + q_{1} q_{2} \qsch_{1423}^P{\left(x; y \right)} \\
&+ q_{2} \left(\left(y_{1} - z_{2}\right) \left(y_{1} - z_{3}\right) + \left(y_{1} - z_{3}\right) \left(y_{5} - z_{1}\right) + \left(y_{5} - z_{1}\right) \left(y_{5} - z_{2}\right)\right) \qsch_{15234}^P{\left(x; y \right)} \\
&+ q_{2} \left(y_{1} + y_{5} + y_{6} - z_{1} - z_{2} - z_{3}\right) \qsch_{162345}^P{\left(x; y \right)} + q_{2} \left(y_{1} + y_{2} + y_{5} - z_{1} - z_{2} - z_{3}\right) \qsch_{25134}^P{\left(x; y \right)} \\
&+ \left(\left(y_{1} - z_{2}\right) \left(y_{1} - z_{3}\right) \left(y_{4} - z_{2}\right) + \left(y_{1} - z_{2}\right) \left(y_{1} - z_{3}\right) \left(y_{5} - z_{1}\right) + \left(y_{1} - z_{2}\right) \left(y_{5} - z_{1}\right) \left(y_{5} - z_{2}\right) \right. \\
&\quad \left. + \left(y_{1} - z_{3}\right) \left(y_{4} - z_{1}\right) \left(y_{5} - z_{1}\right) + \left(y_{4} - z_{1}\right) \left(y_{5} - z_{1}\right) \left(y_{5} - z_{2}\right)\right) \qsch_{15423}^P{\left(x; y \right)} \\
&+ q_{2} \qsch_{1723456}^P{\left(x; y \right)} + q_{2} \qsch_{261345}^P{\left(x; y \right)} + q_{2} \qsch_{35124}^P{\left(x; y \right)} \\
&+ \left(\left(y_{1} - z_{2}\right) \left(y_{1} - z_{3}\right) + \left(y_{1} - z_{3}\right) \left(y_{5} - z_{1}\right) + \left(y_{5} - z_{1}\right) \left(y_{5} - z_{2}\right)\right) \qsch_{156234}^P{\left(x; y \right)} \\
&+ \left(\left(y_{1} - z_{2}\right) \left(y_{1} - z_{3}\right) + \left(y_{1} - z_{2}\right) \left(y_{5} + y_{6} - z_{1} - z_{2}\right) \right. \\
&\quad \left. + \left(y_{1} - z_{3}\right) \left(y_{4} - z_{1}\right) + \left(y_{4} - z_{1}\right) \left(y_{5} + y_{6} - z_{1} - z_{2}\right)\right) \qsch_{164235}^P{\left(x; y \right)} \\
&+ \left(\left(y_{4} - z_{1}\right) \left(y_{5} - z_{3}\right) + \left(y_{4} - z_{1}\right) \left(y_{1} + y_{2} - z_{1} - z_{2}\right) \right. \\
&\quad \left. + \left(y_{5} - z_{2}\right) \left(y_{5} - z_{3}\right) + \left(y_{5} - z_{2}\right) \left(y_{1} + y_{2} - z_{1} - z_{2}\right)\right) \qsch_{25413}^P{\left(x; y \right)} \\
&+ \left(\left(y_{1} - z_{1}\right) \left(y_{4} - z_{1}\right) + \left(y_{5} - z_{2}\right) \left(y_{1} + y_{4} - z_{1} - z_{2}\right)\right) \qsch_{45123}^P{\left(x; y \right)} \\
&+ \left(y_{1} + y_{5} + y_{6} - z_{1} - z_{2} - z_{3}\right) \qsch_{165234}^P{\left(x; y \right)} + \left(y_{1} + y_{2} + y_{5} - z_{1} - z_{2} - z_{3}\right) \qsch_{256134}^P{\left(x; y \right)} \\
&+ \left(y_{1} + y_{4} - z_{1} - z_{2}\right) \qsch_{1742356}^P{\left(x; y \right)} + \left(y_{1} + y_{4} - z_{1} - z_{2}\right) \qsch_{461235}^P{\left(x; y \right)} \\
&+ \left(y_{1} + y_{2} + y_{4} + y_{5} + y_{6} - 2 z_{1} - 2 z_{2} - z_{3}\right) \qsch_{264135}^P{\left(x; y \right)} \\
&+ \left(y_{4} + y_{5} - z_{1} - z_{2}\right) \qsch_{35412}^P{\left(x; y \right)} + \left(y_{4} + y_{5} - z_{1} - z_{2}\right) \qsch_{45213}^P{\left(x; y \right)} \\
&+ \qsch_{1752346}^P{\left(x; y \right)} + \qsch_{265134}^P{\left(x; y \right)} + \qsch_{2741356}^P{\left(x; y \right)} \\
&+ \qsch_{356124}^P{\left(x; y \right)} + \qsch_{364125}^P{\left(x; y \right)} + \qsch_{462135}^P{\left(x; y \right)} \\
&+ \qsch_{561234}^P{\left(x; y \right)}
\end{align*}
For an illustration, consider $d_P = (-,0,1,-,0)$ for the coefficient of $162345$. $q^{d_P} = q_2$. The Peterson-Woodward lift of $d_P$ is $d = (0,0,1,0,0)$, and $w_P = 21354$. $z_P(d) = \{1\}$, so $w_Pw_{z_P(d)} = 12354$. Thus, $ww_Pw_{z_P(d)} = 162435$, and the $q_2$-component of the coefficient of $\qsch_{162345}^P(x;y)$ is $(y_1+y_5+y_6-z_1-z_2-z_3)$, which is the same as the $q_3$-component ($q^d$) of the coefficient of $\qsch_{162435}(x;y)$ in the complete flag variety case (Example \ref{example:cppfpos}).
\end{example}

\begin{corollary}[Triple quantum Schubert positivity for the Grassmannian case]
Suppose $u,v,w\in S_n$ are $k$-Grassmannian, where $P$ is the corresponding parabolic subroup such that $n_1 = k$ and $n_2 = n - k$. Then $c_{u,v}^{w;d}(y;z)$ exhibits triple Schubert positivity.
\end{corollary}


\section{Pieri}

\begin{definition}[$\tom{k}$, $\fx_{k}(u,w)$, $q_k(u,w)$, $E_{p,k}(y;z)$, $\wte(y_A;z)$] \label{definition:pieri}
We define a relation $\tom{k}$ on pairs of permutations, introduced in its non-quantum form by Sottile \cite{sottile}, which was used originally to indicate when the coefficient of a Schubert polynomial is nonzero in a suitable application of the Pieri formula. In the generalization to double Schubert polynomials in \cite{samuelmolev}, the same relation applies to all nonzero coefficients in the expansion of the generalized Pieri formula. The relation plays the same role in our case. 

Given $u,u'\in S_\infty$ and a positive integer $k$, we declare that $u\tom{k} w$ if there exists a $p$ with $0\leq p\leq k$ such that there are transpositions $t_{a_1b_1},\ldots,t_{a_pb_p}$ satisfying 
\begin{enumerate}
\item $a_i\leq k<b_i$ for all $i$
\item $a_i\neq a_j$ whenever $i\neq j$
\item $b_i\geq b_j$ whenever $i<j$
\item We have either
$$\ell(ut_{a_1b_1}\cdots t_{a_ib_i})=\ell(ut_{a_1b_1}\cdots t_{a_{i-1}b_{i-1}})+1$$ 
or 
$$\ell(ut_{a_1b_1}\cdots t_{a_ib_i})=\ell(ut_{a_1b_1}\cdots t_{a_{i-1}b_{i-1}})-2(b_i-a_i)+1$$
for all $i$
\item $w=ut_{a_1b_1}\cdots t_{a_pb_p}$
\end{enumerate}
In contrast to the non-quantum case, $u\tom{k} w$ does not imply that $\ell(u) \leq \ell(w)$, as the final condition can cause the length of the permutation to drop. For intuition, the condition $\ell(ut_{ab}) = \ell(u) - 2(b - a) + 1$ means that if we write the reflection $t_{ab}$ as
$$s_{b}s_{b-1}\cdots s_{a+1}s_as_{a+1}\cdots s_{b-1}s_b$$
then each simple reflection in this word cancels a corresponding simple reflection in a reduced word for $u$, as opposite to the length-increase case where there is a word for which precisely one simple reflection is inserted, regardless of the length of the reflection as a permutation.

If $u\tom{k} w$, define 
$$\fx_{k}(u,w)=\{u(i)\mid i\leq k\mbox{ and }u(i)=w(i)\}$$
and
$$n_k(u,w)=k-\lvert\fx_k(u,w)\rvert.$$
Also, we define $q_k(u,w)$ recursively. If $u=w$, define $q_k(u,w)=1$. If we have defined $q_k(u,w')$ with $w'=t_{a_1b_1}\cdots t_{a_{p-1}b_{p-1}}$, suppose $\ell(w't_{a_pb_p})=\ell(w')+1$. Then define
$$q_k(u,w)=q_k(u,w')$$
If instead $\ell(w't_{a_pb_p})=\ell(w')-2(b_p-a_p)+1$, define
$$q_k(u,w)=q_k(u,w')q_{a_pb_p}$$

Suppose $y$ and $z$ are sets of variables and $0\leq p\leq k$ are integers. We define $E_{p,k}(y;z)$ to be the factorial elementary symmetric polynomial in $y$ and $z$. For specificity, this is the double Schubert polynomial $\sch_{c[p,k]}(y;z)$ where
$$c[p,k] = s_{k-p+1}\cdots s_{k-1}s_k$$
If $p=0$, the polynomial is identically $1$. If $p=k$, $E_{k,k}(y;z)$ can be simply written as
$$E_{k,k}(y;z)=\prod_{i=1}^k(y_i-z_1)$$
An explicit formula can be written for general $p$ without too much trouble as well, using e.g. \cite{molev1999littlewood}:
$$E_{p,k}(y;z)=\sum_{1\leq i_1<\cdots<i_p\leq k}\prod_{j=1}^p(y_{i_j}-z_{i_j-j+1})$$

For a set $A=\{i_1,\ldots,i_{|A|}\}$, we define
$$\wte(y_A;z)=\prod_{i\in A}(y_i-z_1)=E_{|A|,|A|}(y_{i_1},\ldots,y_{i_{|A|}};z),$$
the specialization of the factorial elementary polynomial to which the coefficients of the Pieri formula reduce. It will be used to express those coefficients.
\end{definition}


\begin{theorem}[The Pieri formula] \label{theorem:pieri}
Suppose $u\in S_\infty$ and $k\geq 1$. Then
$$\qsch_u(x;y)\qelem_{k,k}(x;z)=\sum_{u\tom{k} w}q_k(u,w)\wte(y_{\fx_{k}(u,w)};z)\qsch_w(x;y)$$
\end{theorem}



% \begin{theorem} \label{theorem:coeffreduce}
% Suppose $u,v\in S_\infty$ and let $d$ be a weak composition such that $d_i\in \{0,1\}$ for all $i$ and if $d_i=1$, then $d_{i+1}=0$. Let $i_1,\ldots,i_k$ be the integers such that $d_{i_p}=1$ for $1\leq p\leq k$. Then $c_{u,v}^{w,d}(y;z)=0$ unless $\ell(us_{i_p})<\ell(u)$ and $\ell(vs_{i_p})<\ell(v)$ for each $p$, in which case
% $$c_{u,v}^{w,d}(y;z)=c_{us',vs'}^w(y;z)$$
% where
% $$s'=s_{i_1}\cdots s_{i_k}$$
% \end{theorem}

\newcommand{\lom}{\widetilde{\lambda}}
\newcommand{\mom}{\widetilde{\mu}}
\newcommand{\thop}{\widetilde{\theta}}
\newcommand{\medperm}{\overline{\mu}}
\newcommand{\medtheta}{\overline{\theta}}
\newcommand{\Tom}[1]{\xRightarrow{#1}}
\newcommand{\dom}{\mu}

\section{Algorithms for multiplying (quantum) double Schubert polynomials}

\begin{definition}[$\code$, (strict) dominant permutations, $\theta$, $\dom$, $\thop$, $\mom$, medium theta $\medtheta$, $\medperm$] \label{def:codetheta}
We say that a permutation $\mu$ is \emph{dominant} if $\code_i(\mu)\geq \code_{i+1}(\mu)$ for all $i$. A permutation $\mu$ is said to be \emph{strict dominant} if $\mu$ is dominant and whenever $\code_i(\mu)\neq 0$ we have $\code_i(\mu)>\code_{i+1}(\mu)$. For a permutation $v\in S_\infty$, we define a dominant permutation $\dom(v)$ recursively as follows. If $v$ is already dominant, define $\mu(v)=v$. Otherwise, let $i$ be the maximal index such that $\code_i(v)<\code_{i+1}(v)$, and define
$$\dom(v)=\dom(vs_i)$$
Define $\theta_i(v)$ to be $\code_i(\dom(v))$.

Similarly, for a permutation $v\in S_\infty$, we define a strict dominant permutation $\mom(v)$ recursively as follows. If $v$ is already strict dominant, define $\mom(v)=v$. Otherwise, let $i$ be the maximal index such that $\code_{i+1}(v)\neq 0$ and $\code_i(v)\leq \code_{i+1}(v)$, then define
$$\mom(v)=\mom(vs_i)$$
Define $\thop_i(v)$ to be $\code_i(\mom(v))$.

Between these two lies the \emph{medium theta} $\medtheta(v)$, the partition actually iterated by the quantum multiplication algorithm. Like $\theta(v)$ it is a weakly decreasing partition obtained from $\code(v)$ by incrementing and reordering, but it stops short of separating every pair of equal parts: a pair of equal consecutive nonzero parts $\medtheta_i(v)=\medtheta_{i+1}(v)$ is left in place exactly when it is preceded by a strict drop of at least two, that is $\medtheta_{i-1}(v)\geq\medtheta_i(v)+2$ (or $i=1$), and every other ascent or equality is removed by incrementing as for $\theta$. Thus $\medtheta(v)$ permits some, but not all, repeated parts, and lies between $\theta(v)$ (arbitrary repeats) and $\thop(v)$ (no repeats). We write $\medperm(v)$ for the permutation with $\code(\medperm(v))=\medtheta(v)$ and set $|\medtheta(v)|:=\sum_i\medtheta_i(v)=\ell(\medperm(v))$. The number of nonzero parts of $\medtheta(v)$ equals the number of nonzero parts of $\code(v)$.
\end{definition}

\begin{definition}[$\xRightarrow{i}$, $\sigma_i(u,v)$, $D_i(u,v)$]
Let $u,v\in S_\infty$ and $i\geq 1$ be an integer. We define $u\Tom{i} v$ if there exist distinct integers $b_1,b_2,\ldots,b_k$ such that $b_j\neq i$ for all $j$, $\ell(ut_{i,b_1}\cdots t_{i,b_j})=\ell(u)+j$ for all $j$, and 
$$v=ut_{i,b_1}\cdots t_{i,b_k}$$
We define
$$\sigma_i(u,v)=\#\{j\mid b_j<i\}$$
and we define
$$D_i(u,v)=\{u(i)\}\cup\{u(j)\mid u(j)\neq v(j)\}$$
\end{definition}

\begin{theorem}[Formula for ordinary double Schubert polynomials] \label{theorem:orddouble}
Suppose $v\in S_\infty$. Then
$$\sch_v(x;z)=\sum_{(v_0,\ldots,v_k)}\prod_{i=1}^k (-1)^{\sigma_i(v_{i-1},v_i)}E_{\theta_i(v)-\ell(v_{i-1},v_i),\theta_i(v)}(x;z_{D_i(v_{i-1},v_i)})$$
Where $v_{i-1}\xRightarrow{i} v_i$ for all $i$, $v_0=1$, and $v_k=v\dom(v^{-1})$.
\end{theorem}


\begin{theorem}[Formula for quantum double Schubert polynomials] \label{theorem:quandouble}
Suppose $v\in S_\infty$. Then
$$\qsch_v(x;z)=\sum_{(v_0,\ldots,v_k)}\prod_{i=1}^k (-1)^{\sigma_i(v_{i-1},v_i)}E_{\thop_i(v)-\ell(v_{i-1},v_i),\thop_i(v)}^q(x;z_{D_i(v_{i-1},v_i)})$$
Where $v_{i-1}\xRightarrow{i} v_i$ for all $i$, $v_0=1$, and $v_k=v\mom(v^{-1})$.
\end{theorem}

Coupled with the Pieri formulas, these formulas give an algorithm for computing products of (quantum) double Schubert polynomials. We describe its classical form by building it from two ingredients: the elementary Pieri step that generates the $u$-side covers, and the signed path expansion on the $v$-side.

\begin{definition}[$k$-strip cover, admissible $k$-chain, the $(p,k)$-Pieri problem] \label{def:pieriproblem}
Let $w\in S_n$, let $1\leq k<n$, and let $\beta\in\{1,\ldots,n\}\cup\{\infty\}$. A \emph{$k$-strip cover of $(w,\beta)$} is a transposition $t_{ab}$ of positions $a,b$ with $1\leq a\leq k<b\leq n$ such that
\begin{enumerate}
\item $w\lessdot wt_{ab}$ in Bruhat order; that is, $w(a)<w(b)$ and no $c$ with $a<c<b$ satisfies $w(a)<w(c)<w(b)$; and
\item $w(b)<\beta$.
\end{enumerate}
Its \emph{output} is the pair $(wt_{ab},\,w(b))$. An \emph{admissible $k$-chain of length $p$} from $u$ is a sequence $u=w_0,w_1,\ldots,w_p$ carrying bounds $\infty=\beta_0>\beta_1>\cdots>\beta_p$ in which $(w_t,\beta_t)$ is the output of a $k$-strip cover of $(w_{t-1},\beta_{t-1})$ for every $t$. The \emph{$(p,k)$-Pieri problem} for $u$ asks to enumerate every admissible $k$-chain from $u$ of length at most $p$. These chains index the permutations occurring in $e_p(x_1,\ldots,x_k)\,\sch_u$, where $e_p$ is the elementary symmetric polynomial.
\end{definition}

\begin{method}[$\textbf{\textsc{pieri}}(u,p,k)$] \label{method:pieri}
Given $u\in S_n$, $1\leq k<n$, and $0\leq p\leq k$, the algorithm returns a list $T$ of pairs $(w,t)$, one for each admissible $k$-chain from $u$ of length $t\leq p$, recording its endpoint $w$.
\begin{enumerate}
\item Initialize $T\gets[(u,0)]$ and a frontier $F\gets[(u,\infty)]$.
\item For $t=1,\ldots,p$, form a new frontier $F'\gets[\,]$ and, for each $(w,\beta)\in F$:
  \begin{enumerate}
  \item compute $A\gets\{\,a\in\{1,\ldots,k\}\mid w(a)<\beta\,\}$;
  \item for each $b=k+1,\ldots,n$ with $w(b)<\beta$, and each $a\in A$ for which $t_{ab}$ is a $k$-strip cover of $(w,\beta)$, append $(wt_{ab},\,w(b))$ to $F'$ and $(wt_{ab},\,t)$ to $T$.
  \end{enumerate}
  Then set $F\gets F'$.
\item Return $T$.
\end{enumerate}
\end{method}

\begin{lemma}[Multiplicity-freeness] \label{lemma:pierimultfree}
For each $t$, distinct admissible $k$-chains from $u$ of length $t$ have distinct endpoints; equivalently, the endpoint map is a bijection from admissible $k$-chains of length $t$ onto the permutations $w$ with $w\gtrdot^{k} u$ of length $\ell(u)+t$ produced by the Pieri rule for $e_t(x_1,\ldots,x_k)$.
\end{lemma}
\begin{proof}
The strictly decreasing bounds $\beta_0>\beta_1>\cdots$ force the values moved leftward across position $k$ to be strictly decreasing, which is exactly the condition making the $e_t$-Pieri expansion multiplicity-free; see \cite{bergeron2002skew}. Reconstructing the chain from its endpoint by peeling off these moves in decreasing order of value shows the map is injective, hence bijective onto the Pieri support.
\end{proof}

\iffalse
\begin{proposition}[Complexity of $\textbf{\textsc{pieri}}$] \label{prop:piericomplexity}
Let $M:=|T|$ be the number of admissible $k$-chains from $u$ of length at most $p$ that $\textbf{\textsc{pieri}}(u,p,k)$ returns. The algorithm runs in
$$O\!\left(k\,(n-k)\,n\cdot M\right)$$
elementary operations, that is $O(n^3)$ per enumerated permutation. Moreover, by Lemma~\ref{lemma:pierimultfree},
$$M\;\leq\;\#\{\,w\in S_n\mid w\geq u,\ \ell(w)\leq \ell(u)+p\,\}.$$
\end{proposition}
\begin{proof}
Processing one frontier node $(w,\beta)$ costs $O(k)$ to build $A$, and then iterates over the $n-k$ choices of $b$ and the at most $k$ elements of $A$; for each pair the Bruhat-cover test and the transposition each scan $O(n)$ positions. Thus a node costs $O\!\big(k+(n-k)\,k\,n\big)=O\!\big(k\,(n-k)\,n\big)$. By Lemma~\ref{lemma:pierimultfree} the number of nodes ever placed on a frontier equals $\sum_{t=0}^{p-1}N_t$, where $N_t$ is the number of length-$t$ chains; since $\sum_{t=0}^{p-1}N_t\leq\sum_{t=0}^{p}N_t=M$, the total cost is $O\!\big(k\,(n-k)\,n\cdot M\big)$. The bound on $M$ holds because every endpoint satisfies $w\geq u$ and $\ell(w)\leq\ell(u)+p$.
\end{proof}
\fi

We now describe the $v$-side primitive. Its starting point is the structure of a \emph{dominant} double Schubert polynomial. If $\mu$ is dominant with $\code(\mu)=(\lambda_1,\ldots,\lambda_m)$, then $\sch_\mu(x;z)$ factors as a product of factorial elementary symmetric polynomials of maximal degree, one per column of the shape,
$$\sch_\mu(x;z)=\prod_{i=1}^{m}E_{\lambda_i,\,\lambda_i}(x;z^{(i)}),$$
where $z^{(i)}$ is the shifted $z$-alphabet attached to column $i$ (Definition~\ref{definition:pieri}). For a general $v$ one writes $\sch_v(x;z)$ by applying divided differences in the $z$-variables to the dominant polynomial $\sch_{\dom(v^{-1})}(x;z)$; because that polynomial factors by column, the divided-difference action decouples and may be carried out one column (one $z^{(i)}$-group) at a time. Applied to a single column this produces a \emph{signed} descending expansion. The following primitive carries it out.

\begin{definition}[$k$-strip descent beneath an edge, the down-Pieri problem] \label{def:kdownproblem}
Let $w,\mu\in S_n$ and $1\leq k<n$. Say $w'$ lies \emph{beneath the edge} $\mu$ if
$$\mathrm{inv}(w'\mu)=\mathrm{inv}(\mu)-\mathrm{inv}(w'),$$
that is, $\ell(w'\mu)=\ell(\mu)-\ell(w')$. A \emph{$k$-strip descent} of $w$ is a Bruhat descent $w\gtrdot wt_{ab}$ realized by a transposition $t_{ab}$ that moves the value in the pivot position $k$ either leftward (from a position $a<k$, with sign $-1$) or rightward (to a position $b>k$, with sign $+1$), touching only positions still holding their original value. Given a starting label $w$, the edge $\mu$, a degree $p$, and the column $k$, the \emph{down-Pieri problem} asks for the signed list of those length-$t$ descending $k$-strip chains from $w$ ($t\leq p$) whose endpoint lies beneath $\mu$.
\end{definition}

\begin{method}[$\textbf{\textsc{kdown}}(w,\mu,p,k)$] \label{method:kdown}
The algorithm returns a signed list $D$ of triples $(w',t,s)$, one for each length-$t$ descending $k$-strip chain from $w$ ($t\leq p$) whose endpoint $w'$ lies beneath $\mu$, with $s\in\{+1,-1\}$ its sign.
\begin{enumerate}
\item Initialize $D\gets[\,]$; if $w$ is already beneath $\mu$, append $(w,0,+1)$. Set a frontier $F\gets[(w,+1)]$.
\item For $t=1,\ldots,p$, form $F'\gets[\,]$ and, for each $(w'',s)\in F$:
  \begin{enumerate}
  \item (leftward) for each position $a<k$ with $w''(a)=w(a)$ such that $w''\gtrdot w''t_{a,k}$ is a Bruhat descent, set $w'''\gets w''t_{a,k}$ and append $(w''',-s)$ to $F'$; if $w'''$ is beneath $\mu$, append $(w''',t,-s)$ to $D$;
  \item (rightward) for each position $b>k$ with $w''(b)=w(b)$ such that $w''\gtrdot w''t_{k,b}$ is a Bruhat descent, set $w'''\gets w''t_{k,b}$ and append $(w''',+s)$ to $F'$; if $w'''$ is beneath $\mu$, append $(w''',t,+s)$ to $D$.
  \end{enumerate}
  Then set $F\gets F'$.
\item Return $D$.
\end{enumerate}
\end{method}

Applied without the edge $\mu$, the signed descending expansion produces many extraneous terms that ultimately cancel. The edge is exactly the bound that discards them: only endpoints beneath $\mu$ are retained in $D$. In the assembled $v$-side (Method~\ref{method:vpaths} below), only these retained labels seed the next column, so each stage passes forward only labels compatible with the remaining shape.

\begin{method}[$\textbf{\textsc{vpaths}}(v)$, the $v$-side path dictionaries] \label{method:vpaths}
Let $\theta:=\theta(v^{-1})$ with nonzero parts in columns $1,\ldots,L$, and let $vmu:=v\,\dom(v^{-1})$. Processing columns from $i=L$ down to $i=1$, maintain for each column a set of live labels (initialized at column $L$ to the single label $vmu$). At column $i$, with $k=i$ and edge $\mu_i:=\dom$ of the shape truncated below column $i$, call $\textbf{\textsc{kdown}}(w,\mu_i,\theta_i,k)$ for each live label $w$; record its signed output triples as the transitions out of $w$, and let their endpoints be the live labels seeding column $i-1$. The collected transitions form the path dictionaries used by $\textbf{\textsc{schubmult}}$.
\end{method}

\begin{lemma}[Correctness of the edge filter] \label{lemma:edgefilter}
A permutation $w'$ produced by a descending $k$-strip chain from $w$ is retained by $\textbf{\textsc{kdown}}(w,\mu,p,k)$ if and only if $w'$ lies beneath $\mu$, i.e. $\mathrm{inv}(w'\mu)=\mathrm{inv}(\mu)-\mathrm{inv}(w')$. In the assembly of Method~\ref{method:vpaths} the retained signed labels are exactly the terms of $\sch_v(x;z)$ contributed through column $i$, so no extraneous (cancelling) label propagates to column $i-1$.
\end{lemma}
\begin{proof}
The retention test is the identity $\mathrm{inv}(w'\mu)=\mathrm{inv}(\mu)-\mathrm{inv}(w')$ checked in step (2); this is precisely the condition that $w'\leq\mu$ with $\ell(w'\mu)=\ell(\mu)-\ell(w')$. The factorization of the dominant polynomial by column, together with the column-wise decoupling of the $z$-divided differences, identifies the beneath-$\mu$ endpoints at column $i$ with the surviving monomials of $\sch_v(x;z)$ after processing that column; the extraneous chains are exactly those leaving the shape, which the edge discards before they seed column $i-1$.
\end{proof}

\iffalse
\begin{proposition}[Complexity of $\textbf{\textsc{kdown}}$] \label{prop:kdowncomplexity}
Let $R:=|D|$ be the number of retained signed triples and $F^\ast$ the total number of frontier nodes generated by $\textbf{\textsc{kdown}}(w,\mu,p,k)$. The algorithm runs in
$$O\!\left(n^2\cdot F^\ast\right)$$
elementary operations, and $R\leq F^\ast$. Consequently $\textbf{\textsc{vpaths}}(v)$, which invokes $\textbf{\textsc{kdown}}$ once per live label across the $L\leq\min(n-1,\ell_v)$ columns and by Lemma~\ref{lemma:edgefilter} propagates only retained labels, runs in $O\!\big(n^2\cdot\Phi_v\big)$ operations, where $\Phi_v$ is the total number of frontier nodes generated over all columns.
\end{proposition}
\begin{proof}
Processing one frontier node scans $O(n)$ candidate positions, each requiring an $O(n)$ Bruhat-descent test and $O(n)$ transposition, so $O(n^2)$ per node; summing over the $F^\ast$ nodes gives $O(n^2F^\ast)$. Each retained triple arises from a distinct frontier node, so $R\leq F^\ast$. In $\textbf{\textsc{vpaths}}$, by Lemma~\ref{lemma:edgefilter} the labels seeding column $i-1$ are only the retained (beneath-edge) endpoints of column $i$, so the live-label count never includes extraneous paths; summing the per-column frontier work gives $O(n^2\Phi_v)$.
\end{proof}
\fi

\begin{method}[The algorithm \textbf{\textsc{schubmult}}] \label{method:schubmult}
Fix $n$ and let $u,v\in S_n$. The algorithm $\textbf{\textsc{schubmult}}(u,v)$ computes the expansion of $\sch_u(x;z)\,\sch_v(x;z)$ in the Schubert basis as follows. Set $\theta:=\theta(v^{-1})$ and let
$$L:=\#\{i\mid\theta_i\neq 0\}$$
be the number of nonzero parts. The algorithm proceeds in $L$ layers. It maintains a \emph{front}, a collection of nodes each consisting of an intermediate permutation on the $u$-side paired with a label on the $v$-side, initialized to the single node $(u,\,e)$. At layer $i$ it replaces each node by the nodes obtained by combining one $u$-side cover, generated by $\textbf{\textsc{pieri}}(\,\cdot\,,\theta_i,\theta_i)$ (Method~\ref{method:pieri}), with one $v$-side transition, generated by $\textbf{\textsc{vpaths}}(v)$ (Method~\ref{method:vpaths}), accumulating coefficients in the coefficient ring. After all $L$ layers the surviving coefficients at $v$-label $v\dom(v^{-1})$ give the structure constants.
\end{method}

\paragraph{Quantum version.}
The same path organization gives the quantum multiplication algorithm. Its
layers are indexed by the medium theta $\medtheta(v^{-1})$; equal consecutive
parts that are retained in this shape are handled together by the paired
quantum Pieri step. On the $u$-side, classical covers are replaced by the
quantum covers of Remark~\ref{remark:qkbruhat}, including length-decreasing
moves, whose contributions carry the corresponding $q$-monomials. The
$v$-side continues to use the descending path dictionaries of
Method~\ref{method:vpaths}. Thus the two algorithms share the same
layer-by-layer structure, while their Pieri transitions reflect the
classical or quantum multiplication rule, respectively.

The algorithmic construction is one way to compute products. We now turn to
a complementary operator description, which packages the quantum
multiplication rules in the Fomin--Gelfand--Postnikov difference operators.

\subsection{Difference operators and multiplication coefficients}

\iffalse
\begin{definition}[$\ell_u$, $\ell_v$, $P_u$, $B$, $\Delta$] \label{def:complexityquantities}
For $u,v\in S_n$ write $\ell_u:=\ell(u)$ and $\ell_v:=\ell(v)$, and set
$$P_u:=\#\{\,w\in S_n\mid w\geq u \text{ and } \ell(w)\leq \ell_u+\ell_v\,\}.$$
Write $\theta:=\theta(v^{-1})$. The two \emph{per-layer branchings} are
$$B:=\max_{i}\ \max_{w}\ \bigl|\textbf{\textsc{pieri}}(w,\theta_i,\theta_i)\bigr|,\qquad
\Delta:=\max_{i}\ \max_{\lambda}\ \#\{\text{kdown triples }(\lambda',c,\varepsilon)\text{ carried by }\lambda\text{ in column }i\}:$$
$B$ is the maximum number of $u$-side covers generated by a single $\textbf{\textsc{pieri}}$ step (Proposition~\ref{prop:piericomplexity}), and $\Delta$ is the maximum number of $v$-side kdown triples generated from a single label (Proposition~\ref{prop:kdowncomplexity}). Each branching enumerates a family of $k$-strips indexed by (nonempty) subsets of at most $n$ eligible positions, so
$$B\leq 2^{n}-1,\qquad \Delta\leq 2^{n}-1;$$
the bound on $B$ is sharp, attained at $w=w_0$ (where $\textbf{\textsc{pieri}}(w_0,n-1,n)$ produces exactly $2^{n}-1$ covers).
\end{definition}

\begin{theorem}[Complexity of \textbf{\textsc{schubmult}}] \label{thm:complexity}
The asymptotic complexity of $\textbf{\textsc{schubmult}}(u,v)$ is
$$O\!\left(\,\min(n-1,\ell_v)\cdot P_u\cdot 2^{\ell_v}\cdot B\cdot\Delta\cdot n\,\right)$$
operations in the coefficient ring, where $B$ and $\Delta$ are the $u$-side (Pieri) and $v$-side (kdown-triple) per-layer branchings of Definition~\ref{def:complexityquantities}. This is $O\!\left(2^{\,\binom{n}{2}+n\log_2 n}\right)$ in the worst case, and $O\!\left(2^{\,n\log_2 n}\right)$ when $v$ is dominant.
\end{theorem}
\begin{remark}
Note that this does not exceed the time complexity of multiplication of two arbitrary polynomials in $n$ variables in the monomial basis of degree at most $\binom{n}{2}$, which is $O\!\left(2^{\,\binom{n}{2}}\right)$.
\end{remark}
\begin{proof}
By Method~\ref{method:schubmult} the algorithm runs in $L=\#\{i\mid\theta_i(v^{-1})\neq 0\}\leq\min(n-1,\ell_v)$ layers. Write $\mu:=\dom(v^{-1})$, so that $\sum_i\theta_i(v^{-1})=\ell(\mu)$; note that $\ell(\mu)\geq\ell_v$, with equality if and only if $v$ is dominant, so the total degree processed on the shape is in general strictly larger than $\ell_v$. The key point is that the two sides are yoked: at layer $i$ a $u$-side degree $u_i$ and a $v$-side degree $v_i$ are paired with $u_i+v_i=\theta_i$, and the $v$-side descends from $v\mu$ to the identity, so $\sum_i v_i=\ell(v\mu)$. Hence the total degree added on the $u$-side is
$$\sum_i u_i=\sum_i\theta_i-\sum_i v_i=\ell(\mu)-\ell(v\mu)=\ell_v,$$
using $\ell(v\mu)=\ell(\mu)-\ell(v)$ (as $v^{-1}\leq\mu$ in weak order). Thus every intermediate $u$-side permutation satisfies $w\geq u$ and $\ell(w)\leq\ell_u+\ell_v$, so there are at most $P_u$ of them. The $v$-side labels lie in the lower interval $[e,v\mu]$; since $\ell(v\mu)=\ell(\mu)$ can exceed $\ell_v$, we bound their number by $2^{\ell(\mu)}$, but the edge filter of Lemma~\ref{lemma:edgefilter} retains only those beneath the running dominant edge, of which there are at most $2^{\ell_v}$ live at once. Hence the front at the start of each layer has at most $P_u\cdot 2^{\ell_v}$ nodes.

We now account for the work \emph{exactly}. Write $\Phi_i$ for the set of front nodes $(w,\lambda)$ present at the start of layer $i$, each a $u$-side permutation $w$ paired with a $v$-side label $\lambda$. The innermost body of the algorithm matches one $\textbf{\textsc{pieri}}$ cover of $w$ against one kdown triple of $\lambda$, and is executed precisely
$$W\;=\;\sum_{i=1}^{L}\ \sum_{(w,\lambda)\in\Phi_i}\beta_i(w)\,\delta_i(\lambda)$$
times, where $\beta_i(w):=\bigl|\textbf{\textsc{pieri}}(w,\theta_i,\theta_i)\bigr|$ is the $u$-side Pieri branching of Proposition~\ref{prop:piericomplexity} and $\delta_i(\lambda)$ is the number of kdown triples of Proposition~\ref{prop:kdowncomplexity} carried by $\lambda$ in column $i$. This is a genuine \emph{product}: at every node the $v$-side kdown branching $\delta_i(\lambda)$ multiplies---rather than adds to---the $u$-side Pieri branching $\beta_i(w)$, so the signed kdown factor enters the running time as a bona fide multiplicative branching, not as a lower-order additive term. Bounding each factor by its maximum, $\beta_i(w)\leq B$ and $\delta_i(\lambda)\leq\Delta$, and the node count by $|\Phi_i|\leq P_u\cdot 2^{\ell_v}$, gives
$$W\;\leq\;L\cdot P_u\cdot 2^{\ell_v}\cdot B\cdot\Delta,$$
and each executed body costs $O(n)$ ring and permutation operations; this is the stated bound.

For the worst case substitute $P_u\leq n!\leq 2^{\,n\log_2 n}$, $\ell_v\leq\binom{n}{2}$, and $B,\Delta\leq 2^{n}-1$. Because $\log_2(B\,\Delta)\leq 2n=O(n)$, the two branchings---in particular the kdown-triple factor $\Delta$---contribute only $O(n)$ to the exponent, which is absorbed by the $n\log_2 n$ term; hence $W=2^{\,\binom{n}{2}+n\log_2 n+O(n)}=O\!\left(2^{\,\binom{n}{2}+n\log_2 n}\right)$, and the leading exponent $\binom{n}{2}$ is unchanged. If $v$ is dominant then $\theta(v)=\code(v)$, $\ell(\mu)=\ell_v$, every path step carries a $+1$ sign, and each layer's $v$-side is a single path: the number of live labels is $O(\ell_v)$ and the kdown branching collapses to $\Delta=1$, so no term is produced and later canceled. The bound becomes $O(\ell_v\cdot P_u\cdot B\cdot n)=O(2^{\,n\log_2 n})$.
\end{proof}

\begin{remark}
The quantum algorithm has the same overall shape---$L$ layers, a front of $(u\text{-side},\,v\text{-side})$ nodes---but three features make it genuinely more involved than the classical case, all confined to the $u$-side; the $v$-side is \emph{identical}.

First, the layer sequence is the \emph{medium theta} $\medtheta(v^{-1})$ of Definition~\ref{def:codetheta}, not the sorted partition $\theta(v^{-1})$ used in the classical algorithm. Classically a repeated part of $\theta(v^{-1})$ costs nothing extra: it amounts to multiplying by a product of two elementary symmetric polynomials in the same variables, which is a single ordinary Pieri step. In the quantum setting there is no such simple description, so the algorithm instead iterates over $\medtheta(v^{-1})$, which is \emph{almost} strict: an equal consecutive pair $\medtheta_i=\medtheta_{i+1}$ is retained only when it is preceded by a strict drop of at least two, and every such pair is processed as a single \emph{paired} step by a quantum double elementary symmetric Pieri rule (the routine \texttt{double\_elem\_sym\_q}), whose $u$-side branching is the \emph{product} of two single-column Pieri branchings rather than one. The number of nonzero parts of $\medtheta(v^{-1})$ is at most $\min(n-1,\ell_v)$---forming $\medtheta$ from $\code(v^{-1})$ creates no new nonzero part---and, since a paired step consumes two parts at once, so is the number of steps.

Second, the $v$-side is unaffected: it uses the same increasing-path (kdown-triple) machinery as the classical Theorem~\ref{theorem:quandouble}, so the number of $v$-labels is still at most $2^{\ell_v}$ and the kdown-triple branching $\Delta$ is unchanged. Only the column shape it runs over changes, from $\theta(v^{-1})$ to $\medtheta(v^{-1})$, and this does not affect either bound.

Third, on the $u$-side the quantum covers of Remark~\ref{remark:qkbruhat} may \emph{decrease} length and carry a $q$-degree. Two consequences follow. (a) Because the $u$-side climbs the full over-shape $\medperm(v^{-1})$ before the $v$-side filters it back down, an intermediate $u$-permutation $w$ can have length \emph{exceeding} $\ell_u+\ell_v$; the correct bound is $\ell(w)\leq\ell_u+|\medtheta(v^{-1})|$, and we simply bound the number of reachable $w$ by $n!$. (b) A node is a pair $(w,d)$ of a permutation with a $q$-degree, and the same $w$ may in principle be reached at several $q$-degrees; grading by $q$-degree keeps the transition graph acyclic, since each length-decreasing cover $t_{ab}$ raises $\sum_i d_i$ by the strictly positive amount $b-a$. In the implementation the coefficient attached to each front node is in fact a \emph{single} $q$-monomial, because the $q$-degree there is determined by the pair (intermediate permutation, $v$-label); the number of distinct $q$-degrees a coefficient could carry is at most $Q_{u,v}$ (Definition~\ref{def:quantumfront}), which we use as a safe bound on coefficient size.
\end{remark}

\begin{definition}[$\widehat P_{u,v}$, $B_q$, $Q_{u,v}$] \label{def:quantumfront}
For $u,v\in S_n$ set

$$\widehat P_{u,v}:=\#\{\,w\in S_n\mid \ell(w)\leq \ell_u+|\medtheta(v^{-1})|\,\}\;\leq\; n!,$$
and let $B_q$ be the maximum number of covers produced by a single-column quantum Pieri step $\textbf{\textsc{pieri}}^q(w,\medtheta_i,\medtheta_i)$ (the quantum analogue of $B$; the routine \texttt{elem\_sym\_perms\_q}). Since each quantum cover is a quantum $k$-strip, indexed by a subset of at most $n$ eligible positions of one of two types (length-increasing or quantum length-decreasing),
$$B_q\leq 4^{n}=2^{O(n)},\qquad\text{and a paired step contributes at most } B_q^2=2^{O(n)}.$$
Finally set
$$Q_{u,v}:=\binom{\lfloor(\ell_u+\ell_v)/2\rfloor+n-1}{n-1},$$
an upper bound on the number of $q$-monomials in any coefficient (a nonzero $q^d$ requires $\sum_i d_i\leq\tfrac12(\ell_u+\ell_v)$).
\end{definition}

\begin{theorem}[Complexity of quantum \textbf{\textsc{schubmult}}] \label{thm:quantumcomplexity}
The asymptotic complexity of the quantum form of $\textbf{\textsc{schubmult}}(u,v)$ is
$$O\!\left(\,\min(n-1,\ell_v)\cdot \widehat P_{u,v}\cdot 2^{\ell_v}\cdot B_q^{2}\cdot\Delta\cdot Q_{u,v}\cdot n\,\right)$$
operations in the coefficient ring, where the $v$-side kdown-triple factor $\Delta$ is identical to the classical case, the per-layer $u$-side branching is at most $B_q$ for a single column and at most $B_q^2$ for a paired column, and $Q_{u,v}$ bounds the size of each $q$-coefficient. In the worst case this is at most
$$2^{\binom{n}{2}}\cdot n!\cdot Q_{u,v}\cdot 2^{O(n)},$$
whose leading exponent is $\binom{n}{2}$, the same as in the classical bound of Theorem \ref{thm:complexity}.
\end{theorem}

\begin{proof}
\emph{Layers.} The quantum algorithm runs over the nonzero parts of $\medtheta(v^{-1})$, processing a retained pair of equal consecutive parts as one paired step; hence the number of steps is at most the number of nonzero parts of $\medtheta(v^{-1})$, which is at most $\min(n-1,\ell_v)$ (forming $\medtheta(v^{-1})$ from $\code(v^{-1})$ creates no new nonzero part, and there are at most $\min(n-1,\ell_v)$ nonzero entries to begin with).

\emph{The $v$-side is unchanged.} As in Theorem~\ref{thm:complexity}, the number of live $v$-labels is at most $2^{\ell_v}$ and the per-label kdown-triple branching is at most $\Delta$; Theorem~\ref{theorem:quandouble} uses the same increasing-path structure, so nothing here differs from the classical case.

\emph{The $u$-side.} The $u$-side climbs the over-shape $\medperm(v^{-1})$, so every intermediate permutation $w$ satisfies $\ell(w)\leq\ell_u+|\medtheta(v^{-1})|$; length-decreasing quantum covers only lower $\ell(w)$, so this bound persists and there are at most $\widehat P_{u,v}\leq n!$ reachable permutations. At a single-column layer the $u$-branching is at most $B_q$; at a paired (equal-part) layer the routine \texttt{double\_elem\_sym\_q} composes two quantum Pieri steps, so the $u$-branching is at most $B_q^2$. Both are $2^{O(n)}$.

\emph{Coefficients.} Each front coefficient is a $\mathbb Z$-combination of monomials $q^d$ with $\sum_i d_i\leq\tfrac12(\ell_u+\ell_v)$, of which there are at most $Q_{u,v}$; hence each coefficient-ring operation costs $O(Q_{u,v}\cdot n)$. (In fact each front coefficient is a single $q$-monomial, the $q$-degree being determined by the node, so this factor is pessimistic; we retain it as a rigorous upper bound.)

\emph{Assembly.} Exactly as in Theorem~\ref{thm:complexity}, the work is the layered product $\sum_i\sum_{(w,\lambda)\in\Phi_i}\beta^q_i(w)\,\delta_i(\lambda)$ with $\beta^q_i\leq B_q^2$ and $\delta_i\leq\Delta$; bounding the node count by $\widehat P_{u,v}\cdot 2^{\ell_v}$ and multiplying by $B_q^2\cdot\Delta$, the $O(Q_{u,v}\cdot n)$ per-operation cost, and the $\min(n-1,\ell_v)$ layers gives the stated bound. For the worst case use $\widehat P_{u,v}\leq n!$, $\ell_v\leq\binom{n}{2}$, and $B_q,\Delta=2^{O(n)}$; since $\log_2\!\big(n!\,B_q^2\,\Delta\,Q_{u,v}\big)=O(n\log n)=o\!\big(\binom{n}{2}\big)$, the leading exponent is $\binom{n}{2}$, exactly as in the classical bound.
\end{proof}
\fi
%\begin{definition}[$\chi_{j;n}$, $\chi^{(n)}$]




\begin{definition}[$\chi_{j;n}$, $\chi^{(n)}$]
Recall the definition of the commuting difference operators $\chi_{j;n}$, with $j\leq n$, defined by Fomin, Gelfand, and Postnikov, as follows:
$$\chi_{j,n}=x_j-\sum_{1\leq i<j}q_{ij}\partial^{t_{ij}}+\sum_{j<k\leq n} q_{jk}\partial^{t_{jk}}$$
We denote the sequence $(\chi_{1;n},\ldots,\chi_{n;n})$ with $n$ fixed by $\chi^{(n)}$. Recall (CITATION) that if $\sch_v(x)$ has at most $n$ variables, then
$$\qsch_v(\chi^{(n)})(1) = \sch_v(x)$$
and hence
$$\qsch_v(\chi^{(n)};z)(1) = \sch_v(x;z)$$
We specify a more precise formula, as follows.
\end{definition}

\begin{theorem}
Suppose $v\in S_n$. Then
$$\qsch_v(\chi^{(n)};z)=\sum_{u\in S_{n},d} q^dc_{u,v}^{1,d}(x;z)\partial^u$$
\end{theorem}
\begin{proof}
Note that if $\sch_u(x)$ has at most $n$ variables, then
$$\qsch_u(\chi^{(n)};y)\qsch_v(\chi^{(n)};z)=\sum_{w,d} q^dc_{u,v}^{w,d}(y;z)\qsch_w(\chi^{(n)};y)$$
We thus have that
$$(\qsch_u(\chi^{(n)};y)\qsch_v(\chi^{(n)};z))(1)=\sum_{w,d} q^dc_{u,v}^{w,d}(y;z)\sch_w(x;y)$$
This is also equal to
$$\qsch_v(\chi^{(n)};z)(\sch_u(x;y))$$
due to the fact that $(\qsch_u(\chi^{(n)};y))(1)=\sch_u(x;y)$. To obtain the coefficient of $\partial^u$ in $\qsch_v(\chi^{(n)};z)$, we set $y=x$ after applying this operator. Since $\sch_w(x;x)=0$ unless $w=1$, we have that
$$\qsch_v(\chi^{(n)};z)(\sch_u(x;y))|_{y=x}=\sum_d q^dc_{u,v}^{1,d}(x;z)$$
and the result follows.
\end{proof}

%\begin{lemma}
%Let $k$ be an integer and let $v\in S_\infty$. Then
%$$\sum_d q^dc_{s_k,v}^{1,d}(x;z) = \sum_{i\leq k < j} q_{ij}\sch_{vt_{ij}}(x;z)$$
%\end{lemma}
%\begin{proof}
%This value is equal to
%$$\qsch_{s_k}(\chi^{(n)};y)(\sch_v(x;z))$$
%then setting $y=x$ for sufficiently large $n$. We also have
%$$\qsch_{s_k}(\chi^{(n)};y) = \sch_{s_k}(x;y)+\sum_{1\leq i\leq k<j\leq n}q_{ij}\partial^{t_{ij}}$$
%Thus
%$$\qsch_{s_k}(\chi^{(n)};y)(\sch_v(x;z))=\sum_w c_{s_k,v}^w(y;z)\sch_w(x;y)+\sum_{1\leq i\leq k<j\leq n}q_{ij}\sch_{vt_{ij}}(x;z)$$
%Setting $y=z$, the terms involving $\sch_w(x;y)$ disappear, and we are left with the desired result.
%\end{proof}

\begin{theorem}
Let $k$ be an integer and let $v,w\in S_\infty$. Then
$$\sum_d q^dc_{s_k,v}^{w,d}(y;z)=c_{s_k,v}^{w}(y;z)+\sum_{i\leq k < j} q_{ij}\sch_{vt_{ij}w^{-1}}(y;z)$$
\end{theorem}
\begin{proof}
We have that the desired value is equal to the coefficient of $\sch_w(x;y)$ in
$$\qsch_{s_k}(\chi^{(n)};y)(\sch_{v}(x;z))$$
We have that
$$\qsch_{s_k}(\chi^{(n)};y)=\sch_{s_k}(x;y)+\sum_{i\leq k<j\leq n}q_{ij}\partial^{t_{ij}}$$
Applying this operator, we obtain
$$\sum_{i\leq k<j\leq n} q_{ij}\sch_{vt_{ij}}(x;z)+\sum_w c_{s_k,v}^w(y;z)\sch_w(x;y)$$
Applying the Cauchy formula, this is equal to
$$\sum_{i\leq k<j,u} q_{ij}\sch_{vt_{ij}u^{-1}}(y;z)\sch_u(x;y)+\sum_w c_{s_k,v}^w(y;z)\sch_w(x;y)$$
Setting $u=w$ and allowing $n$ to be sufficiently large, we obtain the result.
\end{proof}

% \begin{proof}[Proof of Theorem \ref{theorem:coeffreduce}]
% It suffices to prove this for $v=w_0(n)$. We have
% $$\qsch_{w_0(n)}(\chi;z)=\sum_{u,d} q^dc_{u,w_0(n)}^{1,d}(x;z)\partial^u$$
% Fixing the particular $d$ in the statement of the theorem, we must have that $u=s'$. By the Pieri formula,
% $$c_{s',w_0(n)}^{1,d}(x;z)=\sch_{w_0(n)s'}(x;z)$$
% Thus if we compute the coefficient of $q^d\sch_w(x;y)$ in 
% $$(\qsch_{w_0(n)}(\chi;z))(\sch_u(x;y))$$
% we obtain
% $$c_{us',w_0(n)s'}^w(y;z)$$
% and the result follows.
% \end{proof}

\section{Quantum to classical reduction for Grassmannian products}

\begin{theorem}[{\cite[Theorem~2.7]{huang2015equivariant}}] \label{theorem:quantclass}
	We have the following.
	\begin{enumerate}
	\item $c_{u,v}^{w,d}(y;y)=0$ unless $\mathrm{sgn}_i(w)+\langle d\rangle_i\leq \mathrm{sgn}_i(u)+\mathrm{sgn}_i(v)$ for all $i$.
	\item If $\mathrm{sgn}_k(w)+\langle d\rangle_k=\mathrm{sgn}_k(u)+\mathrm{sgn}_k(v)=2$ for some $k$, then
	$$c_{u,v}^{w,d}(y;y)=c_{us_k,vs_k}^{ws_k,d-\delta_k}(y;y)=\begin{cases}
	c_{us_k,v}^{ws_k,d-\delta_k}(y;y) & \text{if $\mathrm{sgn}_k(w)=0$}\\
	c_{us_k,v}^{ws_k,d}(y;y) & \text{if $\mathrm{sgn}_k(w)=1$}
	\end{cases}$$
	\end{enumerate}
	This is symmetric in $u$ and $v$.
\end{theorem}

\begin{corollary}[{\cite[Corollary~2.8]{huang2015equivariant}}] \label{corollary:quantclass28}
	Let $u,v,w\in S_\infty$ and $d\in Q^\vee$. Fix an index $i$.
	\begin{enumerate}
	\item If $\langle d\rangle_i=\mathrm{sgn}_i(u)=0$ and $\mathrm{sgn}_i(w)=\mathrm{sgn}_i(v)=1$, then
	$$c_{u,v}^{w,d}(y;y)=c_{u,vs_i}^{ws_i,d}(y;y).$$
	\item If $\langle d\rangle_i=\mathrm{sgn}_i(u)=1$ and $\mathrm{sgn}_i(w)=\mathrm{sgn}_i(v)=0$, then
	$$c_{u,v}^{w,d}(y;y)=c_{us_i,v}^{ws_i,d-\delta_i}(y;y).$$
	\end{enumerate}
	This is symmetric in $u$ and $v$.
\end{corollary}

\begin{lemma} \label{lemma:quantumtoordinary}
We have that
$$c_{u,v}^{w,d}(y;z)=\sum_{v_1\cdot v_2=v}c_{u,v_2}^{w,d}(y;y)\sch_{v_1}(y;z)$$
\end{lemma}

\begin{theorem}[Quantum reduction holds for triple quantum Schubert coefficients]\label{theorem:quantclassz}
We have the following.
	\begin{enumerate}
	\item $c_{u,v}^{w,d}(y;z)=0$ unless $\mathrm{sgn}_i(w)+\langle d\rangle_i\leq \mathrm{sgn}_i(u)+\mathrm{sgn}_i(v)$ for all $i$.
	\item If $\mathrm{sgn}_k(w)+\langle d\rangle_k=\mathrm{sgn}_k(u)+\mathrm{sgn}_k(v)=2$ for some $k$, then
	$$c_{u,v}^{w,d}(y;z)=c_{us_k,vs_k}^{ws_k,d-\delta_k}(y;z)=\begin{cases}
	c_{us_k,v}^{ws_k,d-\delta_k}(y;z) & \text{if $\mathrm{sgn}_k(w)=0$}\\
	c_{us_k,v}^{ws_k,d}(y;z) & \text{if $\mathrm{sgn}_k(w)=1$}
	\end{cases}$$
	\end{enumerate}
	This is symmetric in $u$ and $v$.
\end{theorem}
\begin{proof}
For (1), we note that since $\mathrm{sgn}_i(v_2)\leq \mathrm{sgn}_i(v)$ for any factorization $v_1\cdot v_2=v$, we have that $\mathrm{sgn}_i(w)+\langle d\rangle_i\leq \mathrm{sgn}_i(u)+\mathrm{sgn}_i(v_2)$ is a necessary condition for $c_{u,v_2}^{w,d}(y;y)\neq 0$. The result then follows from Lemma \ref{lemma:quantumtoordinary}.

For (2), we note that if $\mathrm{sgn}_k(w)+\langle d\rangle_k=\mathrm{sgn}_k(u)+\mathrm{sgn}_k(v)=2$, then since $\mathrm{sgn}_k(v_2)=\mathrm{sgn}_k(v)$ for any factorization $v_1\cdot v_2=v$ such that $c_{u,v_2}^{w,d}(y;y)\neq 0$, Theorem \ref{theorem:quantclass}(2) holds for each term in the sum in Lemma \ref{lemma:quantumtoordinary}, and the result follows.
\end{proof}

\begin{corollary}[$z$-version of Corollary \ref{corollary:quantclass28}(2)] \label{corollary:quantclass28z}
Let $u,v,w\in S_\infty$ and $d\in Q^\vee$, and fix an index $i$ with $\mathrm{sgn}_i(v)=0$. If $\langle d\rangle_i=\mathrm{sgn}_i(u)=1$ and $\mathrm{sgn}_i(w)=0$, then
$$c_{u,v}^{w,d}(y;z)=c_{us_i,v}^{ws_i,d-\delta_i}(y;z).$$
\end{corollary}
\begin{proof}
Since $\mathrm{sgn}_i(v_2)\leq \mathrm{sgn}_i(v)=0$ for every factorization $v_1\cdot v_2=v$, the hypotheses of Corollary \ref{corollary:quantclass28}(2) hold for each term $c_{u,v_2}^{w,d}(y;y)$ appearing in Lemma \ref{lemma:quantumtoordinary}, so $c_{u,v_2}^{w,d}(y;y)=c_{us_i,v_2}^{ws_i,d-\delta_i}(y;y)$. Resumming via Lemma \ref{lemma:quantumtoordinary} gives the claim.
\end{proof}

% \begin{lemma}
% Suppose $c_{u,v}^{w,d}(y;z)\neq 0$ and $0\leq p\leq k$. Then for all $w'$ such that $u\tom{k} w'$ we have that $c_{u,v}^{w',d'}(y;z)\neq 0$ if and only if there exist $d = d_0, d_1,\ldots,d_p = d'$ and
% $w=w_0\tom{k} w_1\tom{k}\cdots\tom{k} w_p=w'$
% such that for each $0< a \leq p$, either
% \begin{enumerate}
% \item $d_a = d_{a-1}$ and $\ell(w_a) = \ell(w_{a-1}) + 1$, or
% \item $\ell(w_a) = \ell(w_{a-1}) - 1$ and $\langle d_a\rangle_k = \langle d_{a-1}\rangle_k + 2$, $\langle d_a\rangle_{k-1} = \langle d_{a-1}\rangle_{k-1} - 1$, $\langle d_a\rangle_{k+1} = \langle d_{a-1}\rangle_{k+1} - 1$, and $\langle d_a\rangle_j = \langle d_{a-1}\rangle_j$ for all other $j$.
% \item There exist $q\leq k \leq r$ with $q\neq r$ such that 
% $$\ell(w_a) = \ell(w_{a-1}) - 2(r-q) + 1$$
% and $\langle d_a\rangle_q = \langle d_{a-1}\rangle_q + 1$, $\langle d_a\rangle_r = \langle d_{a-1}\rangle_r + 1$, $\langle d_a\rangle_{q-1} = \langle d_{a-1}\rangle_{q-1} - 1$, $\langle d_a\rangle_{r+1} = \langle d_{a-1}\rangle_{r+1} - 1$, and $\langle d_a\rangle_j = \langle d_{a-1}\rangle_j$ for all other $j$.
% \end{enumerate}
% \end{lemma}

% Thus if $c_{u,v}^{w,d}(y;z) = p(y;z)\neq 0$, then for each $i$ we have that
% $$\langle d\rangle_i \geq \ell(w) + \deg p - \ell(u) - \ell(v)$$

\begin{lemma}[Reduction to a coroot supported at $k$] \label{lemma:offkreduction}
Let $u,v,w\in S_\infty$ and let $d\in Q^\vee$. Suppose $v$ has precisely one descent, which is $k$. If $c_{u,v}^{w,d}(y;z)\neq 0$, then there exist $u',w'\in S_\infty$ and $d'\in Q^\vee$ with $0\leq d'\leq d$ such that
$$c_{u,v}^{w,d}(y;z)=c_{u',v}^{w',d'}(y;z)$$
and $\langle d'\rangle_i\leq 0$ for all $i\neq k$. Moreover, if $d'\neq 0$ then $\langle d'\rangle_k>0$.
\end{lemma}
\begin{proof}
Suppose $\langle d\rangle_i>0$ for some $i\neq k$, and choose such an $i$ to be maximal. Because $v$ is $k$-Grassmannian, $\mathrm{sgn}_i(v)=0$, so Theorem \ref{theorem:quantclassz}(1) forces
$$\mathrm{sgn}_i(w)+\langle d\rangle_i\leq \mathrm{sgn}_i(u)+\mathrm{sgn}_i(v)=\mathrm{sgn}_i(u)\leq 1.$$
Since $\langle d\rangle_i\geq 1$, we conclude $\langle d\rangle_i=1$, $\mathrm{sgn}_i(u)=1$, and $\mathrm{sgn}_i(w)=0$. These are exactly the hypotheses of Corollary \ref{corollary:quantclass28z} (note $\mathrm{sgn}_i(v)=0$), so
$$c_{u,v}^{w,d}(y;z)=c_{us_i,v}^{ws_i,d-\delta_i}(y;z).$$
Moreover $\langle d\rangle_i=2d_i-d_{i-1}-d_{i+1}=1$ with $d_{i\pm 1}\geq 0$ forces $d_i\geq 1$, so $d-\delta_i\in Q^\vee$ with $0\leq d-\delta_i\leq d$ and $|d-\delta_i|<|d|$. Replacing $(u,w,d)$ by $(us_i,v\text{ fixed},ws_i,d-\delta_i)$ and iterating, the quantity $|d|$ strictly decreases at each step, so the process terminates in a triple $(u',w',d')$ with $\langle d'\rangle_i\leq 0$ for all $i\neq k$.

Finally, for any finitely supported $d'\in Q^\vee$ the values $\langle d'\rangle_i=2d'_i-d'_{i-1}-d'_{i+1}$ telescope to $\sum_i\langle d'\rangle_i=0$. Hence $\langle d'\rangle_k=-\sum_{i\neq k}\langle d'\rangle_i\geq 0$, with equality if and only if $\langle d'\rangle_i=0$ for all $i$, i.e.\ $d'=0$. Thus $d'\neq 0$ implies $\langle d'\rangle_k>0$.
\end{proof}

% \begin{lemma}
% %Suppose $u,v,w$ are such that $u\Tom{i} w$, $v\Tom{i} w$, and $D_i(u,w)=D_i(v,w)$. Then $u=v$.
% Suppose $(v_0,v_1,\ldots,v_n)$ and $(v_0',\ldots,v_n')$ are such that $v_{i-1}\Tom{i} v_i$, $v_{i-1}'\Tom{i} v_i'$ for all $i$, $v_0=v_0'=1$, and $v_n=v_n'=v$ for some $n$. Suppose also that $D_i(v_{i-1},v_i)=D_{i}(v_{i-1}',v_i')$ for all $i$. Then $v_i=v_i'$ for all $i$.
% \end{lemma}
% \begin{proof}
% This is true if $\ell(v_n)=0$. Otherwise, consider the right descent. There are only two places where it can occur.
% \end{proof}


% %Acting on $321$-avoiding, all the elem syms are Schubs.

% \begin{theorem} \label{theorem:nocan}
% No two terms in Theorem \ref{theorem:orddouble} directly cancel each other.
% \end{theorem}
% \begin{proof}

\section{Proof of the Pieri formula} \label{section:proofpieri}

\newcommand{\ntom}[1]{\centernot{\tom{#1}}}
\newcommand{\desc}{\mathrm{desc}}
\newcommand{\lup}{\triangleleft}
\newcommand{\supp}{\mathrm{supp}}
\newcommand{\Lw}{L}
\newcommand{\qw}{\mathbf{q}}

Throughout this section permutations act on positions: $(vt_{ab})(a)=v(b)$, $(vt_{ab})(b)=v(a)$, and $(vt_{ab})(r)=v(r)$ otherwise. A cycle is written $(a_p,\ldots,a_1,b)$ with the meaning $a_p\mapsto a_{p-1}\mapsto\cdots\mapsto a_1\mapsto b\mapsto a_p$; when $b=\max\{a_1,\ldots,a_p,b\}$ we call $b$ the \emph{top} of the cycle and refer to
\begin{equation}\label{eq:pivot}
(a_p,\ldots,a_1,b)=t_{a_1b}\,t_{a_2b}\cdots t_{a_pb}
\end{equation}
as its \emph{pivot factorization}. (Indeed, the right-hand side sends $a_p\mapsto b\mapsto a_{p-1}$, $b\mapsto a_p$ and $a_1\mapsto b$.)

\subsection{The single-variable Monk formula}

\begin{lemma} \label{lemma:singlevariable}
Let $u\in S_\infty$ and let $i,j\geq 1$. Then
$$(x_i-z_j)\qsch_u(x;y)=(y_{u(i)}-z_j)\qsch_u(x;y)+\sum_{\substack{p\neq i\\\ell(ut_{ip})=\ell(u)+1}}\mathrm{sign}(p-i)\qsch_{ut_{ip}}(x;y)$$
$$+\sum_{\substack{p\neq i\\\ell(ut_{ip})=\ell(u)-2|p-i|+1}}\mathrm{sign}(p-i)q_{ip}\qsch_{ut_{ip}}(x;y)$$
where $q_{ip}$ stands for $q_{\min(i,p)\max(i,p)}$.
\end{lemma}
\begin{proof}
The quantum double Schubert polynomials represent the Schubert classes in the torus-equivariant quantum cohomology of the flag manifold \cite{lam2014quantum}, so they satisfy the equivariant quantum Chevalley formula of Mihalcea \cite[Theorem~1]{mihalcea2007equivariant}. In type $A$, with $\qsch_{s_m}(x;y)=\sum_{a\leq m}(x_a-y_a)$, that formula reads
$$\qsch_{s_m}(x;y)\qsch_u(x;y)=\sum_{a\leq m}(y_{u(a)}-y_a)\,\qsch_u(x;y)+\sum_{\substack{a\leq m<b\\ \ell(ut_{ab})=\ell(u)+1}}\qsch_{ut_{ab}}(x;y)+\sum_{\substack{a\leq m<b\\ \ell(ut_{ab})=\ell(u)-2(b-a)+1}}q_{ab}\qsch_{ut_{ab}}(x;y),$$
the diagonal coefficient being the localization of $\qsch_{s_m}(x;y)$ at $u$, and the two sums ranging over the classical and quantum covers of $u$ in the quantum Bruhat graph (for $\alpha=e_a-e_b$ one has $\langle\omega_m,\alpha^\vee\rangle=1$ exactly when $a\leq m<b$, and $\langle 2\rho,\alpha^\vee\rangle=2(b-a)$). By stability we may work in $S_n$ for $n$ large. Now
$$x_i-z_j=\qsch_{s_i}(x;y)-\qsch_{s_{i-1}}(x;y)+y_i-z_j,$$
with the term $\qsch_{s_{i-1}}$ omitted when $i=1$. Subtracting the two Chevalley formulas, the diagonal coefficients combine to $(y_{u(i)}-y_i)+(y_i-z_j)=y_{u(i)}-z_j$, and the reflections $t_{ab}$ with $a<i<b$ occur in both sums with the same coefficient and cancel. What remains are the reflections $t_{ib}$ with $b>i$ (from $\qsch_{s_i}$, with sign $+1$) and $t_{ai}$ with $a<i$ (from $\qsch_{s_{i-1}}$, with sign $-1$), each carrying the factor $1$ or $q_{ab}$ according to whether the cover is classical or quantum. This is the asserted formula.
\end{proof}

\subsection{Covers and tight cycles}

The following classical count is the only fact about lengths that we shall use.

\begin{lemma} \label{lemma:lengthchange}
Let $v\in S_\infty$ and $a<b$. If $v(a)<v(b)$ then
$$\ell(vt_{ab})=\ell(v)+1+2\,\#\{r\mid a<r<b,\ v(a)<v(r)<v(b)\},$$
and if $v(a)>v(b)$ then
$$\ell(vt_{ab})=\ell(v)-1-2\,\#\{r\mid a<r<b,\ v(b)<v(r)<v(a)\}.$$
In particular $\ell(vt_{ab})\geq\ell(v)+1$ if $v(a)<v(b)$, and $\ell(vt_{ab})\geq\ell(v)-(2(b-a)-1)$ always.
\end{lemma}
\begin{proof}
Passing from $v$ to $vt_{ab}$ changes the inversion status of the pair of positions $(a,b)$, and of the pairs $(a,r)$, $(r,b)$ with $a<r<b$ exactly when $v(r)$ lies strictly between $v(a)$ and $v(b)$ (in which case both pairs change); all other pairs are unaffected.
\end{proof}

\begin{definition}[Covers, $\lup$, $\kappa$] \label{def:cover}
For $v\in S_\infty$ and $a<b$ we write $v\lup vt_{ab}$, and say that $t_{ab}$ is a \emph{cover} of $v$, if either $\ell(vt_{ab})=\ell(v)+1$ (a \emph{classical} cover) or $\ell(vt_{ab})=\ell(v)-2(b-a)+1$ (a \emph{quantum} cover). We set $\kappa(v,t_{ab})=1$ in the first case and $\kappa(v,t_{ab})=q_{ab}$ in the second. By Lemma \ref{lemma:lengthchange}, $t_{ab}$ is a classical cover of $v$ if and only if $v(a)<v(b)$ and no $r$ with $a<r<b$ satisfies $v(a)<v(r)<v(b)$, and a quantum cover if and only if $v(a)>v(b)$ and every $r$ with $a<r<b$ satisfies $v(b)<v(r)<v(a)$. These are the covers of the quantum Bruhat graph, and condition (4) of Definition \ref{definition:pieri} says precisely that each step of a witnessing sequence is a cover (Remark \ref{remark:qkbruhat}).
\end{definition}

It is convenient to separate the two values being exchanged from the values lying between them. For distinct numbers $\alpha\neq\beta$ and a finite set $J$ of numbers not containing $\alpha,\beta$, say that the pair $(\alpha,\beta)$ is \emph{$J$-admissible} if either $\alpha<\beta$ and $J\cap(\alpha,\beta)=\emptyset$, or $\alpha>\beta$ and $J\subseteq(\beta,\alpha)$. Thus $t_{ab}$ is a cover of $v$ if and only if $(v(a),v(b))$ is $J$-admissible for $J=\{v(r)\mid a<r<b\}$; the cover is classical when $\alpha<\beta$ and quantum when $\alpha>\beta$.

\begin{lemma} \label{lemma:admissible}
Let $\alpha,\beta,\gamma$ be distinct numbers and $I,J,J'$ finite sets of numbers avoiding them.
\begin{enumerate}
\item[(i)] $(\alpha,\beta)$ is $(J\cup J')$-admissible if and only if it is $J$-admissible and $J'$-admissible. In particular $J'\subseteq J$ and $J$-admissibility imply $J'$-admissibility.
\item[(ii)] (Composition.) If $(\alpha_1,\alpha_0),(\alpha_2,\alpha_1),\ldots,(\alpha_n,\alpha_{n-1})$ are $I$-admissible then so is $(\alpha_n,\alpha_0)$.
\item[(iii)] (Triangle.) If $(\gamma,\beta)$ is $I$-admissible and $(\alpha,\gamma)$ is $(J\cup\{\beta\}\cup I)$-admissible, then $(\alpha,\beta)$ is $I$-admissible and $(\alpha,\gamma)$ is $J$-admissible.
\end{enumerate}
\end{lemma}
\begin{proof}
(i) is immediate from the definition. For (ii), if $I=\emptyset$ there is nothing to prove. Otherwise $I$ cuts the complement of $I$ into intervals (``blocks''), among them $V_-=(-\infty,\min I)$ and $V_+=(\max I,\infty)$. An $I$-admissible pair $(\alpha',\alpha)$ with $\alpha'<\alpha$ has $\alpha,\alpha'$ in the same block, and one with $\alpha'>\alpha$ has $\alpha\in V_-$ and $\alpha'\in V_+$. Hence along the chain the values descend within a block, except for at most one jump from $V_-$ to $V_+$, after which they remain in $V_+$. Either $\alpha_n<\alpha_0$ lie in a common block, or $\alpha_0\in V_-$ and $\alpha_n\in V_+$; in both cases $(\alpha_n,\alpha_0)$ is $I$-admissible. For (iii), the second assertion is (i). For the first there are four cases. If $\gamma<\beta$ and $\alpha<\gamma$: then $\alpha<\beta$ and $I\cap(\alpha,\beta)\subseteq (I\cap(\alpha,\gamma))\cup(I\cap[\gamma,\beta))=\emptyset$. If $\gamma<\beta$ and $\alpha>\gamma$: then $\beta\in(\gamma,\alpha)$ gives $\beta<\alpha$, and $I\subseteq(\gamma,\alpha)$ together with $I\cap(\gamma,\beta)=\emptyset$ gives $I\subseteq(\beta,\alpha)$. If $\gamma>\beta$ and $\alpha<\gamma$: then $\beta\notin(\alpha,\gamma)$ forces $\beta<\alpha$, and $I\subseteq(\beta,\gamma)$ with $I\cap(\alpha,\gamma)=\emptyset$ gives $I\subseteq(\beta,\alpha)$. The case $\gamma>\beta$, $\alpha>\gamma$ cannot occur, since it would require $\beta\in(\gamma,\alpha)$.
\end{proof}

\begin{definition}[Descents of a cycle, $\desc_u(c)$]
Let $c$ be a cycle and suppose $u\in S_\infty$. Write $c$ as
$$(a_1,\ldots,a_n)$$
where $a_n>a_i$ for all $i<n$ (this is a unique way to express the cycle since the $a_j$ are distinct). Then for any $i$ with $1\leq i<n$ we say that $a_i$ is a \emph{descent} of $c$ with respect to $u$ if $u(a_i)>u(a_{i+1})$. We denote the set of all descents of $c$ with respect to $u$ by $\desc_u(c)$.
\end{definition}

In the notation $c=(a_p,\ldots,a_1,b)$ of \eqref{eq:pivot}, with $a_0:=b$, we have $a_j\in\desc_u(c)$ if and only if $u(a_j)>u(a_{j-1})$; the top $b$ is never a descent.

\begin{definition}[$\Lw_u(\sigma)$, $\qw_u(\sigma)$, $m$-shaped, $u$-tight] \label{def:tight}
Let $u,\sigma\in S_\infty$ and let $c_1,\ldots,c_n$ be the nontrivial cycles of $\sigma$, with $c_i=(a_{i,p_i},\ldots,a_{i,1},b_i)$ and $b_i=\max c_i$. Put
$$\Lw_u(\sigma)=\sum_{i=1}^n\Big(p_i-2\sum_{a\in\desc_u(c_i)}(b_i-a)\Big),\qquad \qw_u(\sigma)=\prod_{i=1}^n\prod_{a\in\desc_u(c_i)}q_{ab_i}.$$
Both depend only on the restriction of $u$ to $\supp(\sigma)$, and $\deg\qw_u(\sigma)=\sum_ip_i-\Lw_u(\sigma)$. We say that $\sigma$ is \emph{$m$-shaped} if every nontrivial cycle of $\sigma$ contains exactly one element greater than $m$, and that $\sigma$ is \emph{$u$-tight} if
$$\ell(u\sigma)-\ell(u)=\Lw_u(\sigma).$$
A \emph{pivot chain} of $\sigma$ is a sequence $t_1,\ldots,t_N$ ($N=\sum_ip_i$) of transpositions obtained by interleaving the pivot factorizations \eqref{eq:pivot} of $c_1,\ldots,c_n$ in any manner that preserves the order of the factors within each cycle; since the cycles are disjoint, $t_1\cdots t_N=\sigma$.
\end{definition}

\begin{lemma}[Tightness] \label{lemma:tight}
Let $u,\sigma\in S_\infty$, let $t_1,\ldots,t_N$ be a pivot chain of $\sigma$, and put $u_0=u$, $u_j=u_{j-1}t_j$.
\begin{enumerate}
\item[(i)] If $t_j=t_{a_{i,l}b_i}$ then $u_{j-1}(a_{i,l})=u(a_{i,l})$ and $u_{j-1}(b_i)=u(a_{i,l-1})$ (with $a_{i,0}=b_i$). Thus the $j$-th step exchanges the values $u(a_{i,l})$ and $u(a_{i,l-1})$, and $a_{i,l}\in\desc_u(c_i)$ exactly when the larger of the two sits at the larger position $b_i$ before the step.
\item[(ii)] $\ell(u\sigma)-\ell(u)\geq\Lw_u(\sigma)$.
\item[(iii)] $\sigma$ is $u$-tight if and only if $u_{j-1}\lup u_j$ for every $j$. In that case the $j$-th step is a classical cover if $a_{i,l}\notin\desc_u(c_i)$ and a quantum cover if $a_{i,l}\in\desc_u(c_i)$, and $\prod_j\kappa(u_{j-1},t_j)=\qw_u(\sigma)$.
\item[(iv)] If $\sigma=\rho\tau$ where $\rho$ and $\tau$ are the products of two complementary sets of cycles of $\sigma$, then $\sigma$ is $u$-tight if and only if $\rho$ is $u$-tight and $\tau$ is $u\rho$-tight. Moreover $\Lw_u(\tau)=\Lw_{u\rho}(\tau)$ and $\qw_u(\tau)=\qw_{u\rho}(\tau)$.
\item[(v)] If $\sigma$ is $u$-tight then so is $t_1\cdots t_j$ for every $j$, and $u_j$ is obtained from $u$ by a sequence of covers.
\end{enumerate}
\end{lemma}
\begin{proof}
(i) The position $a_{i,l}$ is moved only by the factor $t_{a_{i,l}b_i}$, so before that factor is applied it carries $u(a_{i,l})$; the position $b_i$ is moved only by the factors of $c_i$, and after $t_{a_{i,l-1}b_i}$ it carries $u(a_{i,l-1})$ (and $u(b_i)$ before any factor of $c_i$). For (ii) and (iii), by (i) and Lemma \ref{lemma:lengthchange} the $j$-th step increases the length by at least $1$ if $a_{i,l}\notin\desc_u(c_i)$ and by at least $-(2(b_i-a_{i,l})-1)$ if $a_{i,l}\in\desc_u(c_i)$; in both cases equality holds exactly when $t_j$ is a cover of $u_{j-1}$, classical in the first case and quantum in the second (a step with $u_{j-1}(a_{i,l})<u_{j-1}(b_i)$ increases the length and so cannot be a quantum cover, and symmetrically). Summing over $j$ gives (ii), and shows that equality in (ii) holds exactly when every step is a cover; the product of the $\kappa$'s is then the product of $q_{a_{i,l}b_i}$ over the descents, which is $\qw_u(\sigma)$. For (iv), $\ell(u\rho\tau)-\ell(u)=(\ell(u\rho)-\ell(u))+(\ell(u\rho\tau)-\ell(u\rho))\geq \Lw_u(\rho)+\Lw_{u\rho}(\tau)=\Lw_u(\rho)+\Lw_u(\tau)=\Lw_u(\sigma)$, where $\Lw_{u\rho}(\tau)=\Lw_u(\tau)$ because $u\rho$ and $u$ agree on $\supp(\tau)$; equality holds if and only if both inequalities are equalities. (v) The partial product $t_1\cdots t_j$ is again a product of disjoint cycles of the form $(a_{i,l},\ldots,a_{i,1},b_i)$ (a cycle of $\sigma$ truncated after its first $l$ factors), and $t_1,\ldots,t_j$ is a pivot chain of it, all of whose steps are covers by (iii).
\end{proof}

With this language the cycle description of $\tom{k}$ becomes a statement in which the chain disappears entirely.

\begin{lemma} \label{lemma:piericycle}
Suppose $u,w\in S_\infty$ and let $k\geq0$ be an integer, and put $\pi=u^{-1}w$. Then $u\tom{k} w$ if and only if $\pi$ is $k$-shaped and $u$-tight. In that case the pivot chain of $\pi$ in which the cycles are taken in order of decreasing tops is a sequence $t_{a_1b_1},\ldots,t_{a_pb_p}$ as in Definition \ref{definition:pieri}, every pivot chain of $\pi$ is a chain of covers, and
$$q_k(u,w)=\qw_u(\pi).$$
In particular $q_k(u,w)$ does not depend on the witnessing sequence, and $\ell(u,w)+\deg q_k(u,w)$ equals the sum of the lengths of the cycles of $\pi$, each diminished by one.
\end{lemma}
\begin{proof}
Suppose first that $t_{a_1b_1},\ldots,t_{a_pb_p}$ satisfy conditions (1)--(5) of Definition \ref{definition:pieri}. The $b_i$ are weakly decreasing, so the factors sharing a common value $b$ of $b_i$ are consecutive, and by \eqref{eq:pivot} their product $t_{a_ib}\cdots t_{a_{i+r}b}$ is the cycle $(a_{i+r},\ldots,a_i,b)$, whose top is $b$ because $a_l\leq k<b$. Cycles belonging to distinct values of $b$ are disjoint, since the $a_i$ are pairwise distinct and no $a_i$ equals any $b_j$. Hence $\pi$ is the product of these cycles, each of which contains exactly one element greater than $k$, namely its top: $\pi$ is $k$-shaped. The given sequence is a pivot chain of $\pi$, and condition (4) says that each of its steps is a cover; so $\pi$ is $u$-tight by Lemma \ref{lemma:tight}(iii), and the recursive definition of $q_k(u,w)$ assigns to each step exactly the factor $\kappa$, whence $q_k(u,w)=\qw_u(\pi)$.

Conversely, suppose $\pi$ is $k$-shaped and $u$-tight. Order the nontrivial cycles of $\pi$ by decreasing tops and take the corresponding pivot chain. The top of each cycle is its unique element exceeding $k$, so condition (1) holds; (2) holds because the cycles are disjoint and the $a$'s within a cycle are distinct; (3) holds by the chosen order; (4) holds by Lemma \ref{lemma:tight}(iii); and (5) is $t_1\cdots t_N=\pi$. The final assertion is $\deg\qw_u(\pi)=\sum_ip_i-\Lw_u(\pi)$.
\end{proof}

\begin{corollary}[Cycle formula for the quantum weight] \label{corollary:cycleweight}
Suppose $u\tom{k}w$, and write the nontrivial cycles in the decomposition of
$u^{-1}w$ as $c_i=(a_{i,p_i},\ldots,a_{i,1},b_i)$, as in Lemma
\ref{lemma:piericycle}. Then
\[
q_k(u,w)=\prod_i\ \prod_{a\in\desc_u(c_i)}q_{a b_i}.
\]
Equivalently, the exponent of $q_r$ in $q_k(u,w)$ is
\[
\#\{(i,a):a\in\desc_u(c_i),\ a\leq r<b_i\}.
\]
Moreover, $q_k(u,w)$ is multiplicative over the cycles of $u^{-1}w$, and the factor attached to a cycle depends only on the values of $u$ on that cycle.
\end{corollary}
\begin{proof}
This is $q_k(u,w)=\qw_u(\pi)$ together with $q_{ab}=q_a\cdots q_{b-1}$.
\end{proof}

We next show that a tight cycle is very constrained: the values of $u$ along it, read cyclically, decrease except at exactly one place. For a cycle $c=(a_p,\ldots,a_1,b)$ with top $b$ and a permutation $u$ we write
$$\beta_0=u(b),\qquad \beta_j=u(a_j)\quad(1\leq j\leq p),$$
so that $a_j\in\desc_u(c)$ if and only if $\beta_j>\beta_{j-1}$, and we put $u_j=u\,t_{a_1b}\cdots t_{a_jb}$. By Lemma \ref{lemma:tight}(i), the $j$-th step of the pivot chain exchanges $\beta_j$ (at $a_j$) with $\beta_{j-1}$ (at $b$); moreover $u_m(a_l)=\beta_{l-1}$ for $l\leq m$ and $u_m(a_l)=\beta_l$ for $l>m$. If $c$ is $u$-tight, Definition \ref{def:cover} gives, for every $j$ and every position $r$ with $a_j<r<b$:
\begin{align}
&\text{if }\beta_j<\beta_{j-1}\text{ (classical step): } u_{j-1}(r)\notin(\beta_j,\beta_{j-1}); \label{eq:Cj}\\
&\text{if }\beta_j>\beta_{j-1}\text{ (quantum step): } u_{j-1}(r)\in(\beta_{j-1},\beta_j). \label{eq:Qj}
\end{align}

\begin{lemma} \label{lemma:onedescent}
Let $c=(a_p,\ldots,a_1,b)$ be a cycle with top $b$ which is $u$-tight. Then $\desc_u(c)$ has at most one element. In particular this applies to every cycle of $u^{-1}w$ whenever $u\tom{k}w$.
\end{lemma}
\begin{proof}
Suppose $i<j$ are two descents (quantum steps) with no descent strictly between them, so that $\beta_i>\beta_{i-1}$, $\beta_j>\beta_{j-1}$ and $\beta_i>\beta_{i+1}>\cdots>\beta_{j-1}$.

If $a_j>a_i$, then \eqref{eq:Qj} for the step $i$ and $r=a_j$ gives $\beta_j=u_{i-1}(a_j)\in(\beta_{i-1},\beta_i)$, so $\beta_{j-1}<\beta_j<\beta_i$. Hence $j\geq i+2$, and there is $m$ with $i<m<j$ and $\beta_m<\beta_j<\beta_{m-1}$. The step $m$ is classical. If $a_j>a_m$, then \eqref{eq:Cj} for the step $m$ and $r=a_j$ is violated by $u_{m-1}(a_j)=\beta_j$. So $a_m>a_j$, and then \eqref{eq:Qj} for the step $j$ and $r=a_m$ requires $u_{j-1}(a_m)=\beta_{m-1}\in(\beta_{j-1},\beta_j)$, contradicting $\beta_{m-1}>\beta_j$.

If $a_j<a_i$, then \eqref{eq:Qj} for the step $j$ and $r=a_i$ gives $u_{j-1}(a_i)=\beta_{i-1}\in(\beta_{j-1},\beta_j)$, so $\beta_{j-1}<\beta_{i-1}<\beta_i$. Hence again $j\geq i+2$ and there is $m$ with $i<m<j$ and $\beta_m<\beta_{i-1}<\beta_{m-1}$. If $a_i>a_m$, \eqref{eq:Cj} for the step $m$ and $r=a_i$ is violated by $u_{m-1}(a_i)=\beta_{i-1}$. So $a_m>a_i$, and \eqref{eq:Qj} for the step $i$ and $r=a_m$ requires $u_{i-1}(a_m)=\beta_m\in(\beta_{i-1},\beta_i)$, contradicting $\beta_m<\beta_{i-1}$.
\end{proof}

\begin{definition}[Redefinition of $\desc_u(c)$]
If $c$ is a tight cycle, we redefine $\desc_u(c)$ to be the unique element of the previously defined set $\desc_u(c)$ if it exists, or otherwise $0$.
\end{definition}

\begin{lemma} \label{lemma:geq}
Let $c=(a_p,\ldots,a_1,b)$ be a cycle with top $b$ which is $u$-tight, and suppose $\desc_u(c)=a_q$ with $q\geq1$. Then
$$\beta_q>\beta_{q+1}>\cdots>\beta_p>\beta_0>\beta_1>\cdots>\beta_{q-1}.$$
In particular $u(a_j)>u(b)$ for all $j\geq q$.
\end{lemma}
\begin{proof}
All steps other than the $q$-th are classical, which gives the two decreasing runs; it remains to prove $\beta_l>\beta_0$ for $l\geq q$.

First let $l=q$. For $l'<q$ with $a_{l'}<a_q$, \eqref{eq:Cj} for the step $l'$ and $r=a_q$ gives $\beta_q=u_{l'-1}(a_q)\notin(\beta_{l'},\beta_{l'-1})$. If some $l'<q$ has $a_{l'}>a_q$, let $l^\ast$ be the least such; then \eqref{eq:Qj} for the step $q$ and $r=a_{l^\ast}$ gives $u_{q-1}(a_{l^\ast})=\beta_{l^\ast-1}\in(\beta_{q-1},\beta_q)$, so $\beta_q>\beta_{l^\ast-1}$, and descending induction on $l'<l^\ast$ using $\beta_q\notin(\beta_{l'},\beta_{l'-1})$ yields $\beta_q>\beta_{l'-1}$ for all $l'\leq l^\ast$, in particular $\beta_q>\beta_0$. If no such $l'$ exists, the same induction starting from $\beta_q>\beta_{q-1}$ gives $\beta_q>\beta_0$.

Now let $l>q$ and assume $\beta_{l-1}>\beta_0$ (and $\beta_l<\beta_{l-1}$). Suppose, for a contradiction, that $\beta_l<\beta_0$. We claim that $a_{l'}<a_l$ and $\beta_l<\beta_{l'}$ for $1\leq l'\leq q-1$. Indeed, assume $\beta_l<\beta_{l'-1}$ (true for $l'=1$). If $a_{l'}>a_l$, then \eqref{eq:Cj} for the step $l$ and $r=a_{l'}$ requires $u_{l-1}(a_{l'})=\beta_{l'-1}\notin(\beta_l,\beta_{l-1})$, which is false as $\beta_l<\beta_{l'-1}\leq\beta_0<\beta_{l-1}$. So $a_{l'}<a_l$, and \eqref{eq:Cj} for the step $l'$ and $r=a_l$ gives $\beta_l\notin(\beta_{l'},\beta_{l'-1})$, whence $\beta_l<\beta_{l'}$. The claim yields $\beta_l<\beta_{q-1}$. Finally consider the quantum step $q$: if $a_q<a_l$, \eqref{eq:Qj} with $r=a_l$ requires $\beta_l\in(\beta_{q-1},\beta_q)$, contradicting $\beta_l<\beta_{q-1}$; if $a_q>a_l$, \eqref{eq:Cj} for the step $l$ with $r=a_q$ requires $u_{l-1}(a_q)=\beta_{q-1}\notin(\beta_l,\beta_{l-1})$, contradicting $\beta_l<\beta_{q-1}<\beta_0<\beta_{l-1}$. Hence $\beta_l>\beta_0$, and the induction is complete.
\end{proof}

\begin{corollary}[Cyclic monotonicity] \label{corollary:cyclic}
Let $c=(a_p,\ldots,a_1,b)$ be a $u$-tight cycle with top $b$. Then the cyclic sequence $(\beta_0,\beta_1,\ldots,\beta_p)$ has exactly one cyclic ascent: either $\beta_0>\beta_1>\cdots>\beta_p$ (and $\desc_u(c)=0$), or $\desc_u(c)=a_q$ with $q\geq1$ and $\beta_q>\cdots>\beta_p>\beta_0>\cdots>\beta_{q-1}$. Consequently
$$\Lw_u(c)=\begin{cases}p&\text{if the ascent is at the wrap }(\beta_p,\beta_0),\\ p-2(b-a_q)&\text{if the ascent is }(\beta_{q-1},\beta_q),\end{cases}\qquad \qw_u(c)=\begin{cases}1,\\ q_{a_qb}.\end{cases}$$
\end{corollary}

We shall repeatedly use the following elementary observation about such sequences. Call a cyclic sequence of distinct numbers \emph{unimodal} if it has exactly one cyclic ascent, i.e.\ if it is a cyclic rotation of a decreasing sequence. Every cyclically consecutive block of a unimodal sequence is unimodal, and so is any subsequence. If $\Gamma$ is unimodal and $\eta\notin\Gamma$, there is exactly one \emph{gap} of $\Gamma$ (a pair of cyclically consecutive entries) into which $\eta$ can be inserted so that the result is again unimodal, namely: the descent gap $(\gamma,\gamma')$ with $\gamma>\eta>\gamma'$ if $\min\Gamma<\eta<\max\Gamma$, and the ascent gap $(\min\Gamma,\max\Gamma)$ if $\eta<\min\Gamma$ or $\eta>\max\Gamma$. We call it the \emph{insertion gap} of $\eta$ in $\Gamma$.

\subsection{The position $k$}

We organize the inductive step of Theorem \ref{theorem:pieri} around the position $k$. The following lemma records two facts about the interaction of $\tom{k}$ and $\tom{k-1}$ at position $k$.

\begin{lemma} \label{lemma:leveldichotomy}
Suppose $u,w\in S_\infty$ and let $k>1$ be an integer.
\begin{enumerate}
\item[(i)] If $u(k)=w(k)$, then $u\tom{k}w$ if and only if $u\tom{k-1}w$, and in that case $q_k(u,w)=q_{k-1}(u,w)$ and $\fx_k(u,w)=\fx_{k-1}(u,w)\cup\{u(k)\}$.
\item[(ii)] If $u(k)\neq w(k)$ and $u\tom{k}w$, then $u\ntom{k-1}w$.
\end{enumerate}
\end{lemma}
\begin{proof}
Let $\pi=u^{-1}w$. If $u(k)=w(k)$ then $k$ is a fixed point of $\pi$, so no nontrivial cycle of $\pi$ contains $k$, and such a cycle has exactly one element greater than $k$ if and only if it has exactly one element greater than $k-1$. Thus $\pi$ is $k$-shaped if and only if it is $(k-1)$-shaped, while tightness and $\qw_u(\pi)$ do not involve $k$ at all; (i) follows from Lemma \ref{lemma:piericycle}. For (ii), the cycle of $\pi$ containing $k$ is nontrivial and, $\pi$ being $k$-shaped, contains an element greater than $k$; so it contains two elements greater than $k-1$ and $\pi$ is not $(k-1)$-shaped.
\end{proof}

For the rest of this section we fix $u,w\in S_\infty$ and $k>1$, and write
$$\pi=u^{-1}w,\qquad C=\text{the cycle of }\pi\text{ containing }k\ (C=\{k\}\text{ if }\pi(k)=k),\qquad B=\max C.$$
For $m\geq0$ we set $c_m(u,w)=q_m(u,w)\wte(y_{\fx_m(u,w)};z)$ if $u\tom{m}w$ and $c_m(u,w)=0$ otherwise. The \emph{Monk summands} at $w$ are indexed by the reflections $t\in\{t_{ak}\mid a<k\}\cup\{t_{kb}\mid b>k\}$: writing $v=wt$ and $\sigma=u^{-1}v=\pi t$, the summand of $t$ is
$$\mathrm{sign}(t)\,\kappa(v,t)\,c_{k-1}(u,v),\qquad \mathrm{sign}(t_{ak})=-1,\ \mathrm{sign}(t_{kb})=+1,$$
where the summand is understood to be $0$ unless $v\lup w$. (These are precisely the terms of \eqref{eq:up} and \eqref{eq:down} below.) The summand of $t$ is therefore nonzero if and only if $v\lup w$ and $u\tom{k-1}v$, that is, $\sigma$ is $(k-1)$-shaped and $u$-tight.

\begin{lemma}[Classification of the Monk neighbors] \label{lemma:neighbors}
Suppose the Monk summand of $t$ is nonzero, and let $P=\fx_{k-1}(u,w)$. Then every nontrivial cycle of $\pi$ not meeting $t$ has exactly one element greater than $k$, and exactly one of the following holds.
\begin{enumerate}
\item[(U)] $t=t_{kB}$ with $B>k$; $C$ contains exactly two elements greater than $k-1$, namely $k$ and $B$; and $\fx_{k-1}(u,v)=P$. Writing $C=(a_p,\ldots,a_1,B)$ with $a_j=k$, the cycles of $\sigma$ are those of $\pi$ other than $C$ together with
$$C'=(a_{j-1},\ldots,a_1,B)\quad\text{and}\quad C''=(a_p,\ldots,a_{j+1},k).$$
\item[(S)] $t=t_{ak}$ with $a\in C$ and $B>k$; $C$ contains exactly two elements greater than $k-1$, namely $k$ and $B$; in the notation of (U), $a=a_m$ for some $m>j$; $\fx_{k-1}(u,v)=P$; and the cycles of $\sigma$ are those of $\pi$ other than $C$ together with
$$C'_m=(a_p,\ldots,a_{m+1},a_m,a_{j-1},\ldots,a_1,B)\quad\text{and}\quad C''_m=(a_{m-1},\ldots,a_{j+1},k).$$
\item[(P)] $t=t_{ak}$ with $a\in C$ and $B=k$; writing $C=(a_p,\ldots,a_1,k)$ we have $a=a_p=\pi(k)=u^{-1}(w(k))$; $\fx_{k-1}(u,v)=P\sqcup\{w(k)\}$; and the cycles of $\sigma$ are those of $\pi$ other than $C$ together with $C_-=(a_{p-1},\ldots,a_1,k)$ (trivial if $p=1$).
\item[(M)] $t=t_{ak}$ with $a\notin C$; then $B=k$, the cycle $D$ of $\pi$ containing $a$ (possibly $D=\{a\}$) satisfies $D\subseteq[1,k-1]$, and $\fx_{k-1}(u,v)=P\setminus\{u(a)\}$, which equals $P$ unless $D=\{a\}$. The cycles of $\sigma$ are those of $\pi$ other than $C$ and $D$ together with the single cycle $E$ with $\supp(E)=\supp(C)\cup\supp(D)$ obtained by setting $E(a)=\pi(k)$, $E(k)=\pi(a)$ and $E=\pi$ elsewhere.
\end{enumerate}
\end{lemma}
\begin{proof}
A cycle of $\pi$ disjoint from $\supp(t)$ is a cycle of $\sigma$, hence has exactly one element greater than $k-1$, and since it does not contain $k$ this element is greater than $k$. Recall that right multiplication of a permutation by a transposition $t_{xy}$ merges the cycles of $x$ and $y$ if they are distinct, and otherwise splits their common cycle $(x_1,\ldots,x_n)$, $x=x_s$, $y=x_t$, $s<t$, into $(x_t,x_{s+1},\ldots,x_{t-1})$ and $(x_s,x_{t+1},\ldots,x_n,x_1,\ldots,x_{s-1})$.

Let $t=t_{kb}$. If $b\notin C$, the cycle of $\sigma$ through $k$ contains $k$ and $b>k$, which is forbidden. So $b\in C$ and $t$ splits $C$ into a cycle through $k$ and a cycle through $b$; each must contain exactly one element greater than $k-1$, so $C\cap(k-1,\infty)=\{k,b\}$ and $b=B$. The explicit form of the two cycles is the splitting formula applied to $(x_1,\ldots,x_{p+1})=(a_p,\ldots,a_1,B)$ with $x_s=k$, $x_t=B$. Since $v$ agrees with $w$ at all positions less than $k$, $\fx_{k-1}(u,v)=P$. This is (U).

Let $t=t_{ak}$ with $a\in C$, so $t$ splits $C$. If $B>k$, the piece not containing $k$ must contain exactly one element greater than $k-1$, which forces $C\cap(k-1,\infty)=\{k,B\}$ and that piece to contain $B$. With $(x_1,\ldots,x_{p+1})=(a_p,\ldots,a_1,B)$, $k=a_j$ and $a=a_m$: if $m<j$ the piece through $a_m$ is $(a_m,a_{j-1},\ldots,a_{m+1})$, which has no element greater than $k-1$ unless it is trivial ($m=j-1$), in which case the other piece $(k,a_{m-1},\ldots,a_1,B,a_p,\ldots,a_{j+1})$ contains both $k$ and $B$; both are forbidden. So $m>j$, and the splitting formula gives $C'_m$ and $C''_m$. Here $v(a_m)=w(k)=u(a_{j-1})\neq u(a_m)$ and $w(a_m)\neq u(a_m)$, so $a_m$ is a fixed position neither for $v$ nor for $w$, and $\fx_{k-1}(u,v)=P$. This is (S). If $B=k$, the piece not containing $k$ has no element greater than $k-1$ and must be trivial, i.e.\ $\sigma(a)=a$, i.e.\ $\pi(k)=a$; in the notation $C=(a_p,\ldots,a_1,k)$ this is $a=a_p$, and $u(a_p)=u(\pi(k))=w(k)$. Then $\sigma$ has the cycle $C_-$, and $v(a_p)=w(k)=u(a_p)$ makes $a_p$ a new fixed position with value $w(k)\notin P$. This is (P).

Finally let $t=t_{ak}$ with $a\notin C$, so $t$ merges $C$ with the cycle $D$ of $a$ into a cycle $E$ containing $k$; $E$ may contain no further element greater than $k-1$, so $C\subseteq[1,k]$, i.e.\ $B=k$, and $D\subseteq[1,k-1]$. We have $v(a)=w(k)=u(\pi(k))\neq u(a)$ because $\pi(k)\in C\not\ni a$, so $a$ is not a fixed position of $v$; it is one of $w$ exactly when $D=\{a\}$. All other positions less than $k$ agree. This is (M).
\end{proof}

\subsubsection*{The case $B>k$: the up-neighbor and the split down-neighbors}

\begin{lemma}[Up/down pairing] \label{lemma:updown}
Let $\psi\in S_\infty$, let $k<B$, let $z_1,\ldots,z_s$ ($s\geq0$) be distinct positions less than $k$, let $X=(z_s,\ldots,z_1,k,B)$ and $w=\psi X$. Consider the following sequences of transpositions, each with product $X$:
\begin{align*}
F_{\mathrm{std}}&:\ t_{kB},\,t_{z_1B},\,\ldots,\,t_{z_sB} &&\text{(the pivot factorization of $X$)},\\
F_{\mathrm{up}}&:\ t_{z_1k},\,\ldots,\,t_{z_sk},\,t_{kB},\\
F_m&:\ t_{z_1k},\,\ldots,\,t_{z_{m-1}k},\,t_{z_mB},\,\ldots,\,t_{z_sB},\,t_{z_mk} &&(1\leq m\leq s).
\end{align*}
Call such a sequence a \emph{cover path} (from $\psi$) if each of its steps is a cover of the permutation reached before it, and define its \emph{weight} to be the product of the $\kappa$'s of its steps.
\begin{enumerate}
\item[(a)] If $F_{\mathrm{std}}$ is a cover path, then $F_{\mathrm{up}}$ is a cover path of the same weight, and no $F_m$ is a cover path.
\item[(b)] If $F_{\mathrm{up}}$ is a cover path but $F_{\mathrm{std}}$ is not, then exactly one $F_m$ is a cover path, and its weight equals that of $F_{\mathrm{up}}$.
\item[(c)] If some $F_m$ is a cover path, then so is $F_{\mathrm{up}}$.
\end{enumerate}
\end{lemma}
\begin{proof}
(The products are $X$: $F_{\mathrm{up}}$ is $(z_s,\ldots,z_1,k)\,t_{kB}$, and $F_m$ is $(z_{m-1},\ldots,z_1,k)\,(z_s,\ldots,z_m,B)\,t_{z_mk}$; both send $k\mapsto B\mapsto z_s\mapsto\cdots\mapsto z_1\mapsto k$.) If $s=0$ then $F_{\mathrm{std}}=F_{\mathrm{up}}$ and there is nothing to prove, so let $s\geq1$. Put
$$\beta=\psi(B),\qquad \zeta_0=\psi(k),\qquad \zeta_i=\psi(z_i)\ (1\leq i\leq s),\qquad I=\{\psi(r)\mid k<r<B\}.$$
The positions strictly between $k$ and $B$ are never moved, so $I$ is the set of values found there at every stage of every path. Let $\phi_m=\psi t_{z_1k}\cdots t_{z_{m-1}k}$ ($1\leq m\leq s+1$); thus $\phi_m(k)=\zeta_{m-1}$, $\phi_m(z_i)=\zeta_{i-1}$ for $i<m$, $\phi_m(z_i)=\zeta_i$ for $i\geq m$, $\phi_m(B)=\beta$, and $\phi_m$ agrees with $\psi$ elsewhere. Let $J_i=\{\phi_i(r)\mid z_i<r<k\}$. We introduce the predicates
\begin{align*}
U_i&:\ (\zeta_i,\zeta_{i-1})\text{ is }J_i\text{-admissible} &&(1\leq i\leq s),\\
T_i&:\ (\zeta_i,\beta)\text{ is }I\text{-admissible} &&(0\leq i\leq s),\\
A_i&:\ (\zeta_i,\zeta_{i-1})\text{ is }(J_i\cup\{\beta\}\cup I)\text{-admissible} &&(1\leq i\leq s).
\end{align*}
$U_i$ says that the $i$-th step of $F_{\mathrm{up}}$ (from $\phi_i$) is a cover, and $T_s$ that its last step is; so $F_{\mathrm{up}}$ is a cover path if and only if $U_1,\ldots,U_s,T_s$ hold. For $1\leq m\leq s+1$ let $\mathcal S_m$ be the statement that $t_{kB},t_{z_mB},\ldots,t_{z_sB}$ is a cover path from $\phi_m$; thus $F_{\mathrm{std}}$ is a cover path if and only if $\mathcal S_1$, and $\mathcal S_{s+1}=T_s$. The first step of $\mathcal S_m$ exchanges $\zeta_{m-1}$ at $k$ with $\beta$ at $B$ across $I$; the step $t_{z_iB}$ ($i\geq m$) exchanges $\zeta_i$ at $z_i$ with $\zeta_{i-1}$ at $B$ across the positions in $(z_i,k)$, which carry the same values as in $\phi_i$ (the positions $z_l$, $l<i$, carry $\zeta_{l-1}$ and the others their $\psi$-values), the position $k$, which carries $\beta$, and the positions in $(k,B)$. Hence
$$\mathcal S_m\iff T_{m-1}\wedge A_m\wedge\cdots\wedge A_s.$$
Let $\mathcal D_m$ ($1\leq m\leq s$) be the statement that $t_{z_mB},\ldots,t_{z_sB},t_{z_mk}$ is a cover path from $\phi_m$; then $F_m$ is a cover path if and only if $U_1,\ldots,U_{m-1}$ and $\mathcal D_m$ hold. The steps of $\mathcal D_m$ are: $D^{(m)}$, exchanging $\zeta_m$ at $z_m$ with $\beta$ at $B$ across $J_m\cup\{\zeta_{m-1}\}\cup I$ (the position $k$ now carries $\zeta_{m-1}$); for $i>m$, $D^{(i)}$, exchanging $\zeta_i$ with $\zeta_{i-1}$ across $J^\ast_i\cup I$, where $J^\ast_i$ consists of the values at the positions of $(z_i,k]$, so that $J^\ast_i=J_i\cup\{\beta\}$ if $z_m>z_i$ and $J^\ast_i=J_i\cup\{\zeta_{m-1}\}$ if $z_m<z_i$ (the position $z_m$ carries $\beta$ and $k$ carries $\zeta_{m-1}$); and $D^{\mathrm{last}}$, exchanging $\beta$ at $z_m$ with $\zeta_{m-1}$ at $k$ across $J^{\ast\ast}=N_m\cup\{\zeta_{l-1}\mid l>m,\ z_l>z_m\}$, where $N_m$ denotes the set of values carried by the positions of $(z_m,k)$ other than $z_{m+1},\ldots,z_s$ (these values are the same in $\phi_m$ and at the end of $\mathcal D_m$), so that $J_m=N_m\cup\{\zeta_l\mid l>m,\ z_l>z_m\}$. For $i>m$ the set across which $D^{(i)}$ is taken contains $J_i$, and the sets of $D^{(m)},\ldots,D^{(s)}$ all contain $\zeta_{m-1}$, the value at $k$.

\emph{Step 1: $\mathcal S_m\Rightarrow\mathcal S_{m+1}$, and $\mathcal S_m$ implies $U_m,\ldots,U_s,T_s$.} By Lemma \ref{lemma:admissible}(iii) applied with $(\alpha,\gamma,\beta)=(\zeta_m,\zeta_{m-1},\beta)$ and $J=J_m$, the conjunction $T_{m-1}\wedge A_m$ implies $T_m\wedge U_m$. The first claim follows, and the second by iterating it down to $\mathcal S_{s+1}=T_s$, using $A_i\Rightarrow U_i$ (Lemma \ref{lemma:admissible}(i)). In particular $\mathcal S_1$ implies that $F_{\mathrm{up}}$ is a cover path.

\emph{Step 2: $\mathcal D_m$ implies $U_m,\ldots,U_s,T_s$.} For $i>m$, $D^{(i)}$ gives $U_i$ by Lemma \ref{lemma:admissible}(i). By the same lemma, $D^{(m)}$ and $D^{(i)}$ ($i>m$) give that $(\zeta_m,\beta)$ and $(\zeta_i,\zeta_{i-1})$ are $I$-admissible, so $(\zeta_s,\beta)$ is $I$-admissible by Lemma \ref{lemma:admissible}(ii): this is $T_s$. Likewise the pairs $(\zeta_m,\beta),(\zeta_{m+1},\zeta_m),\ldots,(\zeta_s,\zeta_{s-1})$ are $\{\zeta_{m-1}\}$-admissible, which as in the proof of Lemma \ref{lemma:admissible}(ii) means:
\begin{equation}\label{eq:twoblocks}
\text{the values }\beta,\zeta_m,\zeta_{m+1},\ldots,\zeta_s\text{ decrease on each side of }\zeta_{m-1}\text{, with at most one jump from below }\zeta_{m-1}\text{ to above it.}
\end{equation}
For $U_m$, admissibility of $(\zeta_m,\beta)$ with respect to $\{\zeta_{m-1}\}$ leaves three possibilities: (1) $\zeta_{m-1}>\beta>\zeta_m$; (2) $\zeta_m>\zeta_{m-1}>\beta$; (3) $\beta>\zeta_m>\zeta_{m-1}$. We use $D^{(m)}$ restricted to $J_m$ and $D^{\mathrm{last}}$ restricted to $N_m$, and treat the remaining elements $\zeta_l$ ($l>m$, $z_l>z_m$) of $J_m$ by \eqref{eq:twoblocks}. In case (1), $J_m\cap(\zeta_m,\beta)=\emptyset$ and $N_m\cap(\beta,\zeta_{m-1})=\emptyset$, while by \eqref{eq:twoblocks} each $\zeta_l$ ($l>m$) is either less than $\zeta_m$ or greater than $\zeta_{m-1}$; hence $J_m\cap(\zeta_m,\zeta_{m-1})=\emptyset$. In case (2), $J_m\subseteq(\beta,\zeta_m)$ and $N_m\cap(\beta,\zeta_{m-1})=\emptyset$ give $N_m\subseteq(\zeta_{m-1},\zeta_m)$, while by \eqref{eq:twoblocks} the jump occurs at $(\beta,\zeta_m)$ and so $\zeta_{m-1}<\zeta_l<\zeta_m$ for $l>m$; hence $J_m\subseteq(\zeta_{m-1},\zeta_m)$. In case (3), $J_m\cap(\zeta_m,\beta)=\emptyset$ and $N_m\subseteq(\zeta_{m-1},\beta)$ give $N_m\subseteq(\zeta_{m-1},\zeta_m)$, while by \eqref{eq:twoblocks} all of $\beta,\zeta_m,\ldots,\zeta_s$ lie above $\zeta_{m-1}$ and decrease, so again $\zeta_{m-1}<\zeta_l<\zeta_m$ for $l>m$; hence $J_m\subseteq(\zeta_{m-1},\zeta_m)$. In each case this is $U_m$. Together with the prefix $U_1,\ldots,U_{m-1}$ this proves (c).

For the remainder assume that $F_{\mathrm{up}}$ is a cover path. Then the cycle $(z_s,\ldots,z_1,k)$ is $\psi$-tight, so by Corollary \ref{corollary:cyclic} the cyclic sequence $\Gamma=(\zeta_0,\zeta_1,\ldots,\zeta_s)$ is unimodal; let $\delta\in\{0,\ldots,s\}$ be the index of its ascent ($\zeta_\delta>\zeta_{\delta-1}$ if $\delta\geq1$; $\delta=0$ means the ascent is the wrap $(\zeta_s,\zeta_0)$), and let $g\in\{0,\ldots,s\}$ be the insertion gap of $\beta$ in $\Gamma$, where gap $g\geq1$ is the gap $(\zeta_{g-1},\zeta_g)$ and gap $0$ is the wrap. Explicitly: if $\min\Gamma<\beta<\max\Gamma$, then $g$ is determined by $\zeta_{g-1}>\beta>\zeta_g$ (indices modulo $s+1$); otherwise $g=\delta$.

\emph{Step 3: $\mathcal S_m$ implies $m>g$.} $\mathcal S_m$ says that the cycle $(z_s,\ldots,z_m,k,B)$ is $\phi_m$-tight, so by Corollary \ref{corollary:cyclic} the cyclic sequence $(\beta,\zeta_{m-1},\ldots,\zeta_s)$ is unimodal; that is, the wrap gap $(\zeta_s,\zeta_{m-1})$ of the block $\Gamma'=(\zeta_{m-1},\ldots,\zeta_s)$ is the insertion gap of $\beta$ in $\Gamma'$. Suppose $1\leq m\leq g$. Then $\zeta_{g-1},\zeta_g$ both lie in $\Gamma'$ and are consecutive there. If $\min\Gamma<\beta<\max\Gamma$ then $\zeta_{g-1}>\beta>\zeta_g$, so the insertion gap of $\beta$ in $\Gamma'$ is $(\zeta_{g-1},\zeta_g)$, which is not the wrap gap of $\Gamma'$ since $g-1\neq s$. If $\beta$ lies outside $[\min\Gamma,\max\Gamma]$ then $g=\delta$, so $\zeta_g=\max\Gamma$ and $\zeta_{g-1}=\min\Gamma$ lie in $\Gamma'$, whence $\beta$ lies outside the range of $\Gamma'$ and its insertion gap in $\Gamma'$ is the ascent gap $(\zeta_{g-1},\zeta_g)$ of $\Gamma'$, again not the wrap gap. In both cases $\mathcal S_m$ fails.

\emph{Step 4: $\mathcal S_{g+1}$ holds; hence $\mathcal S_m$ holds if and only if $m>g$.} We must verify $T_g$ and $A_i$ for $i>g$; by Lemma \ref{lemma:admissible}(i) we treat the sets $J_i$, $\{\beta\}$ and $I$ separately. Admissibility of $(\zeta_i,\zeta_{i-1})$ with respect to $J_i$ is $U_i$. With respect to $\{\beta\}$: if $i>g$ and $i\neq\delta$, the gap $(\zeta_{i-1},\zeta_i)$ is a descent gap and $\beta\notin(\zeta_i,\zeta_{i-1})$, for otherwise it would be the insertion gap of $\beta$; if $i=\delta>g$, then $\beta\in(\min\Gamma,\max\Gamma)=(\zeta_{\delta-1},\zeta_\delta)$, for otherwise $g=\delta$. With respect to $I$, and for $T_g$: the sequence $s_0=\beta,s_1=\zeta_g,\ldots,s_n=\zeta_s$ is a consecutive block of the unimodal sequence obtained by inserting $\beta$ into $\Gamma$, hence unimodal, and $T_s$ says that $(s_n,s_0)$ is $I$-admissible. If $s_0>s_1>\cdots>s_n$, then $T_s$ is classical, $I\cap(s_n,s_0)=\emptyset$, and every pair $(s_i,s_{i-1})$ is $I$-admissible. Otherwise the linear ascent is at some $j\geq1$, so $s_j>\cdots>s_n>s_0>\cdots>s_{j-1}$; then $T_s$ is quantum, $I\subseteq(s_0,s_n)$, and the pairs $(s_i,s_{i-1})$ with $i<j$ lie below $s_0$, the pair $(s_j,s_{j-1})$ has $I\subseteq(s_0,s_n)\subseteq(s_{j-1},s_j)$, and the pairs with $i>j$ lie above $s_n$; all are $I$-admissible. The pairs $(s_i,s_{i-1})$ are exactly $(\zeta_g,\beta)$ and $(\zeta_i,\zeta_{i-1})$, $i>g$. This proves $\mathcal S_{g+1}$, and the second assertion follows from Steps 1 and 3.

\emph{Step 5: $\mathcal D_m$ implies $m=g$.} Restricting $D^{(m)}$ to $\{\zeta_{m-1}\}$ leaves the three cases (1)--(3) of Step 2. In case (1) the gap $(\zeta_{m-1},\zeta_m)$ is a descent gap containing $\beta$, so $g=m$. In case (2) the gap is the ascent gap, $m=\delta$, and $\beta<\zeta_{m-1}=\min\Gamma$, so $g=\delta=m$. In case (3), $m=\delta$ and $\beta>\zeta_m=\max\Gamma$, so again $g=m$.

\emph{Step 6: if $g\geq1$ then $\mathcal D_g$ holds.} Write $a=z_g$. We are in one of the cases (1) $\zeta_{g-1}>\beta>\zeta_g$ ($g\neq\delta$), (2) $g=\delta$ and $\beta<\zeta_{g-1}=\min\Gamma$, (3) $g=\delta$ and $\beta>\zeta_g=\max\Gamma$. For $D^{(g)}$: admissibility of $(\zeta_g,\beta)$ with respect to $I$ is $T_g$ (Step 4); with respect to $\{\zeta_{g-1}\}$ it holds in each of the three cases by inspection; with respect to $J_g$ it follows from $U_g$: in case (1), $J_g\cap(\zeta_g,\zeta_{g-1})=\emptyset$ gives $J_g\cap(\zeta_g,\beta)=\emptyset$; in case (2), $J_g\subseteq(\zeta_{g-1},\zeta_g)\subseteq(\beta,\zeta_g)$; in case (3), $J_g\subseteq(\zeta_{g-1},\zeta_g)$ is disjoint from $(\zeta_g,\beta)$. For $D^{(i)}$, $i>g$: if $z_g>z_i$ the required statement is $A_i$ (Step 4); if $z_g<z_i$ then $z_g\in(z_i,k)$, so $\zeta_{g-1}=\phi_i(z_g)\in J_i$ already, $J^\ast_i=J_i$, and the statement follows from $U_i$ and the $I$-admissibility established in Step 4. For $D^{\mathrm{last}}$, the pair $(\beta,\zeta_{g-1})$ must be admissible with respect to $N_g\subseteq J_g$ and to the values $\zeta_{l-1}$, $l>g$, $z_l>z_g$: in case (1) it is classical, and $(\beta,\zeta_{g-1})\subseteq(\zeta_g,\zeta_{g-1})$ contains no element of $J_g$ (by $U_g$) and no $\zeta_l$ (as $\zeta_{g-1},\zeta_g$ are consecutive in the decreasing order of $\Gamma$); in case (2) it is classical, and $(\beta,\zeta_{g-1})$ lies below $\min\Gamma$ and below $J_g\subseteq(\zeta_{g-1},\zeta_g)$; in case (3) it is quantum, and $J_g\subseteq(\zeta_{g-1},\zeta_g)\subseteq(\zeta_{g-1},\beta)$ while every $\zeta_{l-1}$ with $l>g$ lies in $(\zeta_{g-1},\zeta_g]\subseteq(\zeta_{g-1},\beta)$.

\emph{Step 7: weights.} Suppose $F_{\mathrm{up}}$ is a cover path and $g\geq1$. The two paths $F_{\mathrm{up}}$ and $F_g$ share the prefix $t_{z_1k},\ldots,t_{z_{g-1}k}$; after it, the quantum steps of $F_{\mathrm{up}}$ are $t_{z_\delta k}$ (if $\delta\geq g$) and $t_{kB}$ (if $\zeta_s>\beta$), while those of $F_g$ are $t_{z_gB}$ (if $\zeta_g>\beta$), $t_{z_\delta B}$ (if $\delta>g$) and $t_{z_gk}$ (if $\beta>\zeta_{g-1}$). In case (1) with $\delta>g$ we have $\zeta_s>\zeta_0\geq\zeta_{g-1}>\beta$, and both weights are $q_{z_\delta k}q_{kB}=q_{z_\delta B}$; in case (1) with $\delta<g$ we have $\zeta_s<\zeta_g<\beta$ and both weights are $1$; in case (2) both are $q_{z_gk}q_{kB}=q_{z_gB}$; in case (3) both are $q_{z_gk}$. If instead $g=0$, so that $\mathcal S_1$ holds, the quantum steps of $F_{\mathrm{std}}$ are $t_{kB}$ (if $\zeta_0>\beta$) and $t_{z_\delta B}$ (if $\delta\geq1$); if $\zeta_s>\beta>\zeta_0$ (so $\delta\geq1$) both $F_{\mathrm{std}}$ and $F_{\mathrm{up}}$ have weight $q_{z_\delta B}$; if $\beta<\min\Gamma$ (so $\delta=0$) both have weight $q_{kB}$; and if $\beta>\max\Gamma$ both have weight $1$.

\emph{Conclusion.} (a): $\mathcal S_1$ gives $F_{\mathrm{up}}$ (Step 1) with the same weight (Step 7), and $g=0$ (Step 3), so no $\mathcal D_m$ holds (Step 5). (b): if $F_{\mathrm{up}}$ is a cover path and $\mathcal S_1$ fails, then $g\geq1$ (Step 4), $\mathcal D_g$ holds (Step 6) and no other $\mathcal D_m$ does (Step 5); as $U_1,\ldots,U_{g-1}$ hold, $F_g$ is the unique $F_m$ which is a cover path, and its weight is that of $F_{\mathrm{up}}$ (Step 7). (c) is Step 2.
\end{proof}

\begin{lemma} \label{lemma:pieriknotkn1}
Suppose $u,w\in S_\infty$ and let $k>1$ be an integer. If $u\tom{k} w$ but $u\ntom{k-1} w$, then there is a unique $b>k$ such that $wt_{kb}\lup w$ and $u\tom{k-1} wt_{kb}$, namely $b=B$. Furthermore, in that case there does not exist an $a<k$ such that $wt_{ak}\lup w$ and $u\tom{k-1} wt_{ak}$, and
$$\kappa(wt_{kB},t_{kB})\,q_{k-1}(u,wt_{kB})=q_k(u,w),\qquad \fx_{k-1}(u,wt_{kB})=\fx_k(u,w).$$
\end{lemma}
\begin{proof}
By Lemma \ref{lemma:piericycle}, $\pi$ is $k$-shaped and $u$-tight but not $(k-1)$-shaped, so $u(k)\neq w(k)$ and $B>k$ (Lemma \ref{lemma:leveldichotomy}). Write $C=(a_p,\ldots,a_1,B)$ with $a_j=k$, let $\rho$ be the product of the other nontrivial cycles of $\pi$, let $C'=(a_{j-1},\ldots,a_1,B)$, and put $\psi=u\rho C'$ and $X=(a_p,\ldots,a_{j+1},k,B)$, so that $w=\psi X$ and $X$ has the form of Lemma \ref{lemma:updown} with $s=p-j$ and $z_i=a_{j+i}$. Take the pivot chain of $\pi$ consisting of the pivot factorization of $\rho$, then that of $C$; by Lemma \ref{lemma:tight}(iii),(v) all its steps are covers, $\rho$ is $u$-tight, $C'$ is $u\rho$-tight (it is $C$ truncated after $j-1$ factors), and $F_{\mathrm{std}}$ is a cover path from $\psi$. By Lemma \ref{lemma:updown}(a), $F_{\mathrm{up}}$ is a cover path of the same weight. Its first $s$ steps form the pivot factorization of $C''=(a_p,\ldots,a_{j+1},k)$, so $\rho C'C''=u^{-1}wt_{kB}$ is $u$-tight by Lemma \ref{lemma:tight}(iii),(iv); it is $(k-1)$-shaped, its cycles being those of $\rho$ (each with exactly one element greater than $k$), $C'$ (top $B$, other elements less than $k$) and $C''$ (top $k$, other elements less than $k$). Hence $u\tom{k-1}wt_{kB}$, and the last step of $F_{\mathrm{up}}$ says $wt_{kB}\lup w$. Uniqueness of $b$ is Lemma \ref{lemma:neighbors}(U). The weights are $\qw_u(\rho)\qw_u(C')$ times the weights of $F_{\mathrm{std}}$, respectively $F_{\mathrm{up}}$, which agree; and $\fx_{k-1}(u,wt_{kB})=\fx_{k-1}(u,w)=\fx_k(u,w)$ since $u(k)\neq w(k)$.

Suppose finally that $a<k$ satisfies $wt_{ak}\lup w$ and $u\tom{k-1}wt_{ak}$. By Lemma \ref{lemma:neighbors} we are in case (S): $a=a_m$ with $m>j$, and $\sigma=\rho C'_mC''_m$ is $u$-tight. Taking the pivot chain of $\sigma$ in the order: factors of $\rho$, the first $j-1$ factors of $C'_m$ (which are those of $C'$), the factors of $C''_m$, the remaining factors $t_{a_mB},\ldots,t_{a_pB}$ of $C'_m$, and appending the cover $t_{a_mk}$, we obtain that $F_{m-j}$ is a cover path from $\psi$, contradicting Lemma \ref{lemma:updown}(a).
\end{proof}

\subsubsection*{The case $B=k$: the peel and the merges}

\begin{lemma}[Peel] \label{lemma:peel}
Let $\omega\in S_\infty$ and let $C=(a_p,\ldots,a_1,k)$ ($p\geq1$) be a cycle with top $k$; put $w=\omega C$, $C_-=(a_{p-1},\ldots,a_1,k)$ and $v=wt_{a_pk}=\omega C_-$. Then $C$ is $\omega$-tight if and only if $C_-$ is $\omega$-tight and $v\lup w$, and in that case $\kappa(v,t_{a_pk})\,\qw_\omega(C_-)=\qw_\omega(C)$.
\end{lemma}
\begin{proof}
The pivot factorization of $C$ is that of $C_-$ followed by $t_{a_pk}$, and $\omega C_-=v$. Apply Lemma \ref{lemma:tight}(iii).
\end{proof}

\begin{lemma}[Merge] \label{lemma:merge}
Let $\omega\in S_\infty$, $k>1$, $\eta=\omega(k)$, and let $D$ be a cycle with $\supp(D)\subseteq[1,k-1]$, allowing $D$ to be trivial, $D=\{a\}$ with $a\leq k-1$. Put $w=\omega D$, so $w(k)=\eta$. For $a\in\supp(D)$ let $\tau_a=Dt_{ak}$, the cycle with support $\supp(D)\cup\{k\}$ sending $a\mapsto k\mapsto D(a)$, and $v_a=wt_{ak}=\omega\tau_a$. Then:
\begin{enumerate}
\item[(a)] there is at most one $a\in\supp(D)$ such that $\tau_a$ is $\omega$-tight and $v_a\lup w$;
\item[(b)] such an $a$ exists if and only if $D$ is $\omega$-tight and $\max\supp(D)=k-1$ (for trivial $D=\{a\}$: if and only if $a=k-1$);
\item[(c)] for this $a$, $\kappa(v_a,t_{ak})\,\qw_\omega(\tau_a)=q_{k-1}\,\qw_\omega(D)$.
\end{enumerate}
\end{lemma}
\begin{proof}
\emph{Trivial $D=\{a\}$.} Then $w=\omega$, $\tau_a=t_{ak}$, and $\tau_a$ is $\omega$-tight if and only if $w\lup wt_{ak}$. If also $wt_{ak}\lup w$, then both $\ell(wt_{ak})-\ell(w)$ and its negative lie in $\{1,\,1-2(k-a)\}$, which forces $k-a=1$. Conversely for $a=k-1$ the transposition $t_{k-1,k}$ changes the length by $\pm1$, so both relations hold, one classical and one quantum with $\kappa=q_{k-1}$; thus $\kappa(v_a,t_{ak})\qw_\omega(\tau_a)=q_{k-1}=q_{k-1}\qw_\omega(D)$.

\emph{Nontrivial $D=(d_r,\ldots,d_1,d)$, $r\geq1$, $d=\max\supp(D)$.} Put $d_0=d$, $\gamma_0=\omega(d)$, $\gamma_i=\omega(d_i)$ and $\epsilon=k-d\geq1$. For $a=d_m$ ($0\leq m\leq r$) the cycle $\tau_m:=\tau_{d_m}$ is
$$\tau_m=(d_{m-1},\ldots,d_1,d,d_r,\ldots,d_m,k),$$
whose pivot factorization $t_{d_mk}t_{d_{m+1}k}\cdots t_{d_rk}t_{dk}t_{d_1k}\cdots t_{d_{m-1}k}$ exchanges, in order, the values $\eta,\gamma_m,\gamma_{m+1},\ldots,\gamma_r,\gamma_0,\gamma_1,\ldots,\gamma_{m-1}$: the sequence of values of $\tau_m$ in the sense of Corollary \ref{corollary:cyclic} is the cyclic sequence $\Gamma=(\gamma_0,\ldots,\gamma_r)$ of $D$ with $\eta$ inserted into the gap $(\gamma_{m-1},\gamma_m)$ (the wrap gap $(\gamma_r,\gamma_0)$ when $m=0$). Note $w(d_m)=\gamma_{m-1}$ for $m\geq1$ and $w(d)=\gamma_r$; we write $w(a)$ for the value at $a=d_m$ in either case.

\emph{Necessity.} Suppose $\tau_m$ is $\omega$-tight and $v_m:=v_{d_m}\lup w$. By Corollary \ref{corollary:cyclic} the cyclic sequence of $\tau_m$ is unimodal; deleting $\eta$ shows that $\Gamma$ is unimodal, and then $m$ must be the insertion gap of $\eta$ in $\Gamma$. This already proves (a). Let $\delta$ be the ascent index of $\Gamma$ ($\gamma_\delta>\gamma_{\delta-1}$, or $\delta=0$ for the wrap), so that by Corollary \ref{corollary:cyclic} and Definition \ref{def:tight}
$$\Lw_\omega(D)=r-2(d-d_\delta),\qquad \qw_\omega(D)=q_{d_\delta d}$$
(read as $r$ and $1$ when $\delta=0$). Now compare lengths. Tightness gives $\ell(v_m)-\ell(\omega)=\Lw_\omega(\tau_m)$, Lemma \ref{lemma:tight}(ii) gives $\ell(w)-\ell(\omega)\geq\Lw_\omega(D)$ with equality if and only if $D$ is $\omega$-tight, and $v_m\lup w$ gives $\ell(v_m)-\ell(w)=-1$ if $w(a)>\eta$ and $\ell(v_m)-\ell(w)=2(k-a)-1$ if $w(a)<\eta$. Hence
\begin{equation}\label{eq:mergeineq}
\Lw_\omega(\tau_m)-\Lw_\omega(D)\ \geq\ \begin{cases}-1&\text{if }w(a)>\eta,\\ 2(k-a)-1&\text{if }w(a)<\eta,\end{cases}
\end{equation}
with equality if and only if $D$ is $\omega$-tight. We compute the left side from Corollary \ref{corollary:cyclic} in the three possible configurations of the insertion gap.
\begin{itemize}
\item[(I)] $\min\Gamma<\eta<\max\Gamma$, so $\gamma_{m-1}>\eta>\gamma_m$ and $m\neq\delta$. The ascent of the sequence of $\tau_m$ is the (interior) step onto $\gamma_\delta$, whose position is $d_\delta$, so $\Lw_\omega(\tau_m)=(r+1)-2(k-d_\delta)$ and $\Lw_\omega(\tau_m)-\Lw_\omega(D)=1-2\epsilon$. Here $w(a)=\gamma_{m-1}>\eta$, so \eqref{eq:mergeineq} reads $1-2\epsilon\geq-1$, forcing $\epsilon=1$ and equality.
\item[(II)] $\eta<\min\Gamma$, so $m=\delta$. The ascent of the sequence of $\tau_m$ is its first step $(\eta,\gamma_\delta)$, at position $d_\delta$; again $\Lw_\omega(\tau_m)-\Lw_\omega(D)=1-2\epsilon$, and $w(a)=\gamma_{\delta-1}>\eta$ forces $\epsilon=1$ and equality.
\item[(III)] $\eta>\max\Gamma$, so $m=\delta$. The sequence of $\tau_m$ is decreasing, its ascent being the wrap, so $\Lw_\omega(\tau_m)=r+1$ and $\Lw_\omega(\tau_m)-\Lw_\omega(D)=1+2(d-d_\delta)$. Here $w(a)=\gamma_{\delta-1}<\eta$, and \eqref{eq:mergeineq} reads $1+2(d-d_\delta)\geq2(k-d_\delta)-1$, i.e.\ $\epsilon\leq1$; so $\epsilon=1$ and equality.
\end{itemize}
In every case $d=k-1$ and $D$ is $\omega$-tight, proving the ``only if'' part of (b).

\emph{Sufficiency.} Suppose $D$ is $\omega$-tight and $d=k-1$; then $\Gamma$ is unimodal (Corollary \ref{corollary:cyclic}). Let $m$ be the insertion gap of $\eta$ and $a=d_m$. The computations above show that \eqref{eq:mergeineq} holds with equality, and since $\ell(w)-\ell(\omega)=\Lw_\omega(D)$, the tightness of $\tau_m$ is equivalent to $\ell(v_m)-\ell(w)$ taking the value $-1$ (if $w(a)>\eta$) or $2(k-a)-1$ (if $w(a)<\eta$), i.e.\ to $v_m\lup w$. It therefore suffices to show $wt_{ak}\lup w$, that is, that $(w(a),\eta)$ is $J$-admissible for $J=\{w(r)\mid a<r<k\}$. The positions in $(a,k)$ are: $d=k-1$ (if $a\neq d$), carrying $\gamma_r$; the $d_l>a$, carrying $\gamma_{l-1}$; and positions outside $\supp(D)$, carrying their $\omega$-values. If $a=d=k-1$ the set $J$ is empty and we are done; so assume $m\geq1$.
\begin{itemize}
\item[(I)] $\gamma_{m-1}>\eta>\gamma_m$: the pair $(w(a),\eta)=(\gamma_{m-1},\eta)$ is classical and we need $J\cap(\eta,\gamma_{m-1})=\emptyset$. No $\gamma_l$ lies in $(\gamma_m,\gamma_{m-1})$, these being consecutive in the decreasing order of $\Gamma$. The $m$-th step of the pivot chain of $D$ is classical, so by \eqref{eq:Cj} every position $r\in(d_m,d)$ outside $\supp(D)$ has $\omega(r)\notin(\gamma_m,\gamma_{m-1})$.
\item[(II)] $m=\delta$, $\eta<\gamma_{\delta-1}=\min\Gamma$: the pair $(\gamma_{\delta-1},\eta)$ is classical and we need $J\cap(\eta,\gamma_{\delta-1})=\emptyset$; all $\gamma_l\geq\gamma_{\delta-1}$, and by \eqref{eq:Qj} for the quantum step $\delta$ of $D$ the positions $r\in(d_\delta,d)$ outside $\supp(D)$ have $\omega(r)\in(\gamma_{\delta-1},\gamma_\delta)$.
\item[(III)] $m=\delta$, $\eta>\gamma_\delta=\max\Gamma$: the pair $(\gamma_{\delta-1},\eta)$ is quantum and we need $J\subseteq(\gamma_{\delta-1},\eta)$; the values $\gamma_l$ occurring in $J$ are different from $\gamma_{\delta-1}$ (which sits at $a$) and lie in $(\gamma_{\delta-1},\gamma_\delta]$, and the remaining positions have $\omega(r)\in(\gamma_{\delta-1},\gamma_\delta)$ by \eqref{eq:Qj}.
\end{itemize}
This proves the ``if'' part of (b).

\emph{Weights.} In cases (I) and (II), $\kappa(v_a,t_{ak})=1$ and $\qw_\omega(\tau_m)=q_{d_\delta k}=q_{d_\delta,k-1}q_{k-1}=\qw_\omega(D)q_{k-1}$ (read $q_{dk}=q_{k-1}$ and $\qw_\omega(D)=1$ when $\delta=0$). In case (III), $\kappa(v_a,t_{ak})=q_{ak}=q_{d_\delta k}$ and $\qw_\omega(\tau_m)=1$, giving the same product. This is (c).
\end{proof}

\subsubsection*{The cancellations}

We now assemble the Monk summands. Throughout, $\rho$ denotes the product of the nontrivial cycles of $\pi$ other than $C$ and other than the cycle $D$ of $\pi$ containing $k-1$ (we set $D=\{k-1\}$ if $k-1$ is fixed by $\pi$, and we do not single out $D$ when $k-1\in C$). By Lemma \ref{lemma:neighbors}, Lemma \ref{lemma:piericycle} and Lemma \ref{lemma:tight}(iv), each of the quantities $c_{k-1}(u,w)$, $c_{k-2}(u,w)$, $c_k(u,w)$ and each Monk summand vanishes unless $\rho$ is $u$-tight and each cycle of $\rho$ contains exactly one element greater than $k$; indeed each requires $\rho$ to be a tight product of cycles with exactly one element greater than $k$ or than $k-1$ or than $k-2$, and these conditions agree for cycles containing neither $k$ nor $k-1$. (For $c_{k-2}(u,w)$, note that $k-1$ lies in $C$ or in $D$, hence not in a cycle of $\rho$.) \emph{We assume this from now on}; otherwise all terms considered below are zero.

\begin{lemma}[Reflection--drop cancellation] \label{lemma:reflectioncancel}
Suppose $u,w\in S_\infty$, let $k>1$, and suppose $u(k)=w(k)$. Then
$$\sum_{\substack{a<k,\ wt_{ak}\lup w}}\!\!\!\mathrm{sign}(a-k)\,\kappa_{ak}\,c_{k-1}(u,wt_{ak})\ +\ \sum_{\substack{b>k,\ wt_{kb}\lup w}}\!\!\!\mathrm{sign}(b-k)\,\kappa_{kb}\,c_{k-1}(u,wt_{kb})\ +\ q_{k-1}\,c_{k-2}(u,w)=0,$$
where $\kappa_{rs}=\kappa(wt_{rs},t_{rs})$. In fact the two reflection sums contribute at most one nonzero summand in total: it occurs precisely when $u\tom{k-2}w$ (equivalently, when the quantum term is nonzero), it is indexed by some $a<k$, and it equals $-\,q_{k-1}c_{k-2}(u,w)$.
\end{lemma}
\begin{proof}
Here $C=\{k\}$, so by Lemma \ref{lemma:neighbors} only summands of type (M) can be nonzero: $t=t_{ak}$ with the $\pi$-cycle $D_a$ of $a$ contained in $[1,k-1]$. If $D_a$ is nontrivial then it is not a cycle of $\rho$ (those have an element greater than $k$), so $D_a=D$ and $\max D=k-1$; otherwise $a$ is a fixed point of $\pi$. Let $a$ be such a candidate and put $\omega_a=wD_a^{-1}$, so that $w=\omega_aD_a$, $\omega_a(k)=w(k)$, and $\omega_a=u\rho$ when $D_a=D$. Then $u^{-1}wt_{ak}=(\pi D_a^{-1})\,\tau_a$ with $\tau_a=D_at_{ak}$ as in Lemma \ref{lemma:merge}, where $\pi D_a^{-1}$ is the product of the cycles of $\pi$ other than $D_a$; by Lemma \ref{lemma:tight}(iv), $u^{-1}wt_{ak}$ is $u$-tight if and only if $\pi D_a^{-1}$ is $u$-tight and $\tau_a$ is $\omega_a$-tight. Now apply Lemma \ref{lemma:merge} to $(\omega_a,D_a)$. If $a$ is a fixed point, the summand can be nonzero only for $a=k-1$, which requires $k-1$ to be fixed, i.e.\ $D=\{k-1\}$; in that case $\pi D_a^{-1}=\rho$ and the summand is nonzero. If $D_a=D$ is nontrivial, at most one $a\in D$ works, and one does if and only if $D$ is $u\rho$-tight. Altogether at most one $a$ gives a nonzero summand, and one exists if and only if either $D=\{k-1\}$, or $D$ is nontrivial with $\max D=k-1$ and $u\rho$-tight. On the other hand $u\tom{k-2}w$ if and only if $\pi=\rho D$ is $(k-2)$-shaped and $u$-tight, i.e.\ (given our standing assumption on $\rho$) if and only if $D$ is trivial or has $k-1$ as its unique element greater than $k-2$, and $D$ is $u\rho$-tight: the same condition. When it holds, with $v=wt_{ak}$ the unique contributing neighbor and $\omega=\omega_a=u\rho$, Lemma \ref{lemma:merge}(c) and Lemma \ref{lemma:piericycle} give
$$\kappa_{ak}\,q_{k-1}(u,v)=\qw_u(\rho)\,\kappa_{ak}\qw_\omega(\tau_a)=\qw_u(\rho)\,q_{k-1}\qw_\omega(D)=q_{k-1}\,q_{k-2}(u,w),$$
and Lemma \ref{lemma:neighbors}(M) gives $\fx_{k-1}(u,v)=\fx_{k-1}(u,w)\setminus\{u(a)\}$, which equals $\fx_{k-2}(u,w)$: if $D=\{k-1\}$ then $a=k-1$ and we remove the fixed value $u(k-1)$; if $D$ is nontrivial then neither $a$ nor $k-1$ is a fixed position of $w$. Hence the summand is $-\kappa_{ak}c_{k-1}(u,v)=-q_{k-1}c_{k-2}(u,w)$, cancelling the quantum term. If the condition fails, both the reflection sums and the quantum term vanish.
\end{proof}

\begin{lemma}[Moved reflection cancellation] \label{lemma:movedcancel}
Let $u,w\in S_\infty$, let $k>1$, and suppose $u(k)\neq w(k)$, so that $C$ is nontrivial.
\begin{enumerate}
\item[(a)] If $B>k$, then $c_{k-1}(u,w)=c_{k-2}(u,w)=0$, and the sum of the Monk summands equals $c_k(u,w)$. More precisely: if $u\tom{k}w$ the only nonzero Monk summand is that of $t_{kB}$, and it equals $c_k(u,w)$; if $u\ntom{k}w$, the nonzero Monk summands are either none or exactly two, one of type (U) and one of type (S), with equal absolute value and opposite signs.
\item[(b)] If $B=k$, then $c_k(u,w)=0$, no summand of type (U) occurs, and
$$(y_{w(k)}-z_1)\,c_{k-1}(u,w)\;-\!\!\sum_{\substack{a<k,\ wt_{ak}\lup w}}\!\!\kappa_{ak}\,c_{k-1}(u,wt_{ak})\;+\;q_{k-1}c_{k-2}(u,w)=0.$$
In fact the diagonal term cancels against the summand of type (P), and the quantum term against the (at most one) summand of type (M).
\end{enumerate}
\end{lemma}
\begin{proof}
(a) $C$ contains $k$ and $B$, two elements greater than $k-1$, so $\pi$ is neither $(k-1)$- nor $(k-2)$-shaped and $c_{k-1}(u,w)=c_{k-2}(u,w)=0$. By Lemma \ref{lemma:neighbors}, a nonzero Monk summand is of type (U) or (S), and requires $C\cap(k-1,\infty)=\{k,B\}$; together with the standing assumption on $\rho$ this means that $\pi$ is $k$-shaped, which we may assume (otherwise all Monk summands vanish and $c_k(u,w)=0$). Write $C=(a_p,\ldots,a_1,B)$ with $a_j=k$, $C'=(a_{j-1},\ldots,a_1,B)$, $\psi=u\rho C'$ and $X=(a_p,\ldots,a_{j+1},k,B)$, so that $w=\psi X$ as in Lemma \ref{lemma:updown} with $s=p-j$, $z_i=a_{j+i}$. If $C'$ is not $u\rho$-tight then nothing is tight (Lemma \ref{lemma:tight}(iv),(v): $C'$ is a truncation of $C$ and of each $C'_m$, and a cycle of $u^{-1}wt_{kB}$), all Monk summands vanish, and $c_k(u,w)=0$. So assume $C'$ is $u\rho$-tight. Then, exactly as in the proof of Lemma \ref{lemma:pieriknotkn1}: $u\tom{k}w$ if and only if $F_{\mathrm{std}}$ is a cover path from $\psi$; the summand of type (U) is nonzero if and only if $F_{\mathrm{up}}$ is a cover path, and then equals $\qw_u(\rho)\qw_u(C')\,\mathrm{wt}(F_{\mathrm{up}})\,\wte(y_P;z)$ with $P=\fx_{k-1}(u,w)$; and the summand of type (S) indexed by $a_m$ ($m>j$) is nonzero if and only if $F_{m-j}$ is a cover path, and then equals $-\qw_u(\rho)\qw_u(C')\,\mathrm{wt}(F_{m-j})\,\wte(y_P;z)$ (Lemma \ref{lemma:neighbors}(S) for the fixed set). If $u\tom{k}w$, Lemma \ref{lemma:updown}(a) and Lemma \ref{lemma:pieriknotkn1} show that the only nonzero summand is that of $t_{kB}$ and that it equals $q_k(u,w)\wte(y_{\fx_k(u,w)};z)=c_k(u,w)$. If $u\ntom{k}w$, Lemma \ref{lemma:updown}(b),(c) show that either no summand is nonzero, or exactly two are, namely (U) and one (S), with equal weights; they cancel, and $c_k(u,w)=0$.

(b) $C\subseteq[1,k]$ has no element greater than $k$, so $\pi$ is not $k$-shaped and $c_k(u,w)=0$; type (U) is excluded by Lemma \ref{lemma:neighbors}. Write $C=(a_p,\ldots,a_1,k)$. For the diagonal and peel terms, let $\omega_0=u\rho D$ if $k-1\notin C$ and $\omega_0=u\rho$ if $k-1\in C$ (where $D$ is the cycle of $k-1$, trivial or not), so that $w=\omega_0C$ and $\omega_0$ agrees with $u$ on $\supp(C)$. The diagonal term is nonzero if and only if $\pi$ is $(k-1)$-shaped and tight, i.e.\ (given the standing assumption) $D$ is trivial or has an element greater than $k$, $D$ is $u\rho$-tight, and $C$ is $\omega_0$-tight (the conditions on $D$ being vacuous when $k-1\in C$). The summand of type (P) is nonzero if and only if $\rho D C_-$ is $(k-1)$-shaped and tight and $v=wt_{a_pk}\lup w$, i.e.\ the same conditions on $D$ together with: $C_-$ is $\omega_0$-tight and $v\lup w$. By Lemma \ref{lemma:peel} these two sets of conditions are equivalent, and when they hold, $\kappa_{a_pk}q_{k-1}(u,v)=\qw_u(\rho D)\kappa_{a_pk}\qw_{\omega_0}(C_-)=\qw_u(\rho D)\qw_{\omega_0}(C)=q_{k-1}(u,w)$, while $\fx_{k-1}(u,v)=\fx_{k-1}(u,w)\sqcup\{w(k)\}$ by Lemma \ref{lemma:neighbors}(P). Since $(y_{w(k)}-z_1)\wte(y_{\fx_{k-1}(u,w)};z)=\wte(y_{\fx_{k-1}(u,w)\cup\{w(k)\}};z)$, the diagonal term and the (negatively signed) summand of type (P) cancel.

For the merge and quantum terms, a summand of type (M) requires $a\notin C$ with the cycle $D_a$ of $a$ contained in $[1,k-1]$; as in the proof of Lemma \ref{lemma:reflectioncancel}, either $a$ is fixed by $\pi$ or $D_a=D$ with $\max D=k-1$ (in particular $k-1\notin C$). If $k-1\in C$, then $k-1$ is not fixed and no such $D$ exists; since a fixed-point candidate must be $a=k-1$ by the trivial case of Lemma \ref{lemma:merge} (applied as below), there are no nonzero summands of type (M), and $c_{k-2}(u,w)=0$ since $C$ contains the two elements $k-1,k$ greater than $k-2$. So suppose $k-1\notin C$. For a candidate $a$ put $\omega_a=wD_a^{-1}$, so that $w=\omega_aD_a$ and $\omega_a(k)=w(k)$; then $u^{-1}wt_{ak}=\rho_a\,E$, where $\rho_a$ is the product of the cycles of $\pi$ other than $C$ and $D_a$ and $E=C\tau_a$ is the cycle of Lemma \ref{lemma:neighbors}(M), $\tau_a=D_at_{ak}$. The pivot factorization of $E$ is that of $C$ followed by that of $\tau_a$, so by Lemma \ref{lemma:tight}(iii),(iv), $u^{-1}wt_{ak}$ is $u$-tight if and only if $\rho_a$ is $u$-tight, $C$ is $u\rho_a$-tight and $\tau_a$ is $\omega_a$-tight (note $u\rho_aC=\omega_a$), and then $\qw_u(E)=\qw_u(C)\qw_{\omega_a}(\tau_a)$. By Lemma \ref{lemma:merge} applied to $(\omega_a,D_a)$: a fixed-point candidate contributes only if $a=k-1$, i.e.\ $D=\{k-1\}$; and if $D_a=D$ is nontrivial, at most one $a\in D$ contributes, and one does if and only if $D$ is $\omega_a$-tight, where now $\omega_a=u\rho C=:\omega$. Altogether at most one summand of type (M) is nonzero, and one is if and only if $C$ is $u\rho$-tight and either $D=\{k-1\}$, or $D$ is nontrivial with $\max D=k-1$ and $\omega$-tight. This is exactly the condition for $u\tom{k-2}w$: $\pi=\rho CD$ is $(k-2)$-shaped (as $k-1\notin C$, $C$ has the unique element $k$ greater than $k-2$, and $D$ the unique element $k-1$) and it is $u$-tight if and only if $C$ is $u\rho$-tight and $D$ is $\omega$-tight. When it holds, Lemma \ref{lemma:merge}(c) gives
$$\kappa_{ak}q_{k-1}(u,wt_{ak})=\qw_u(\rho)\qw_u(C)\,\kappa_{ak}\qw_\omega(\tau_a)=\qw_u(\rho)\qw_u(C)\,q_{k-1}\qw_\omega(D)=q_{k-1}q_{k-2}(u,w),$$
and $\fx_{k-1}(u,wt_{ak})=\fx_{k-1}(u,w)\setminus\{u(a)\}=\fx_{k-2}(u,w)$ as in the proof of Lemma \ref{lemma:reflectioncancel}. Thus the summand of type (M) equals $-q_{k-1}c_{k-2}(u,w)$ and cancels the quantum term; when the condition fails, both vanish.
\end{proof}

\begin{lemma}[Vanishing] \label{lemma:vanishing}
Suppose $u,w\in S_\infty$, let $k>1$, and suppose $u\ntom{k}w$. Then
\begin{align*}
(y_{w(k)}-z_1)\,c_{k-1}(u,w)
&+\sum_{\substack{b>k,\ wt_{kb}\lup w}}\kappa_{kb}\,c_{k-1}(u,wt_{kb})
-\sum_{\substack{a<k,\ wt_{ak}\lup w}}\kappa_{ak}\,c_{k-1}(u,wt_{ak})\\
&+q_{k-1}c_{k-2}(u,w)=0,
\end{align*}
where $\kappa_{rs}=\kappa(wt_{rs},t_{rs})$.
\end{lemma}
\begin{proof}
If $u(k)=w(k)$, then $u\ntom{k-1}w$ by Lemma \ref{lemma:leveldichotomy}(i), so the diagonal term vanishes, and the rest is Lemma \ref{lemma:reflectioncancel}. If $u(k)\neq w(k)$ and $B>k$, the diagonal and quantum terms vanish and the Monk summands sum to $c_k(u,w)=0$ by Lemma \ref{lemma:movedcancel}(a). If $u(k)\neq w(k)$ and $B=k$, this is Lemma \ref{lemma:movedcancel}(b).
\end{proof}

We now record the recurrence satisfied by the quantum factorial elementary polynomials, which drives the induction.

\begin{lemma} \label{lemma:elemrecursion}
For all integers $k$ we have
$$\qelem_{k,k}(x;z)=(x_k-z_1)\qelem_{k-1,k-1}(x;z)+q_{k-1}\qelem_{k-2,k-2}(x;z),$$
with the conventions $\qelem_{0,0}=1$ and $\qelem_{k,k}=0$ if $k<0$. Moreover, the ordinary polynomials $\wte(y_A;z)=\prod_{i\in A}(y_i-z_1)$ satisfy, for any set $A$ and any $c\notin A$, the single-variable identity
$$\wte(y_{A\cup\{c\}};z)=(y_c-z_1)\,\wte(y_A;z).$$
\end{lemma}
\begin{proof}
Here $\qelem_{k,k}(x;z)=\sum_{j=0}^k(-z_1)^{k-j}e^{(q)}_j(x_1,\ldots,x_k)=\Delta_k(-z_1\mid x_1,\ldots,x_k)$, and expanding the tridiagonal determinant $\Delta_k(t\mid x_1,\ldots,x_k)$ along its last row gives $\Delta_k=(x_k+t)\Delta_{k-1}+q_{k-1}\Delta_{k-2}$. The second identity is clear.
\end{proof}

The first recurrence is the quantization of the classical relation $E_{k,k}=(x_k-z_1)E_{k-1,k-1}$; the extra term $q_{k-1}\qelem_{k-2,k-2}$ is the only feature absent from the non-quantum setting. Discarding all quantum terms recovers the classical Pieri formula \cite[Theorem~7.1]{samuelmolev}, whose term-by-term cancellation is carried out in \cite[Proof of Proposition~4.4]{samuelmolev} by means of the reflection-matching Lemmas 7.6--7.9 of that paper; the lemmas above are the quantum counterparts, the essential new points being the length-decreasing covers permitted by Lemma \ref{lemma:piericycle}, the attendant bookkeeping of the powers of $q$, and the extra term of the recurrence, which is absorbed by the merge summands of Lemma \ref{lemma:merge}.

\begin{proof}[Proof of Theorem \ref{theorem:pieri}]
For a permutation $w$, write $[\qsch_w]F$ for the coefficient of $\qsch_w(x;y)$ in an expansion $F$ of a quantum double Schubert polynomial in the basis $\{\qsch_w(x;y)\}$. With this notation the theorem is the assertion that, for all $u$ and all $k\geq1$,
$$[\qsch_w]\big(\qsch_u(x;y)\qelem_{k,k}(x;z)\big)=\begin{cases}q_k(u,w)\,\wte(y_{\fx_{k}(u,w)};z) & \text{if }u\tom{k} w,\\[2pt] 0 & \text{otherwise.}\end{cases}$$
For brevity, when $u\tom{m} w$ we set
$$c_m(u,w)=q_m(u,w)\,\wte(y_{\fx_{m}(u,w)};z),$$
and $c_m(u,w)=0$ if $u\ntom{m}w$; note that $u\tom{0}w$ only for $w=u$, with $c_0(u,u)=1$, in accordance with $\qelem_{0,0}=1$. The claim is that $[\qsch_w](\qsch_u\qelem_{k,k})=c_k(u,w)$.

We argue by strong induction on $k$.

\emph{Base case $k=1$.} We have $\qelem_{1,1}(x;z)=x_1-z_1$, and Lemma \ref{lemma:singlevariable} with $i=j=1$ gives
$$(x_1-z_1)\qsch_u(x;y)=(y_{u(1)}-z_1)\qsch_u(x;y)+\sum_{\substack{b>1\\\ell(ut_{1b})=\ell(u)+1}}\qsch_{ut_{1b}}(x;y)+\sum_{\substack{b>1\\\ell(ut_{1b})=\ell(u)-2(b-1)+1}}q_{1b}\qsch_{ut_{1b}}(x;y).$$
The permutations $w$ with $u\tom{1}w$ are exactly $w=u$ and $w=ut_{1b}$ for $b>1$ with $u\lup ut_{1b}$, since Definition \ref{definition:pieri} permits at most one reflection when $k=1$. For $w=u$ we have $q_1=1$, $\fx_{1}(u,u)=\{u(1)\}$, and $\wte(y_{\{u(1)\}};z)=y_{u(1)}-z_1$, matching the diagonal term. For $w=ut_{1b}$ we have $\fx_{1}=\emptyset$ and $\wte(y_\emptyset;z)=1$; if the step is length-increasing then $q_1(u,w)=1$, matching the second sum, and if it is length-decreasing then $q_1(u,w)=q_{1b}$, matching the third sum. All other coefficients vanish on both sides. This proves the base case.

\emph{Inductive step.} Fix $k>1$ and assume the claim at all levels $m<k$. By Lemma \ref{lemma:elemrecursion},
\begin{equation}\label{eq:pieriexpand}
\qsch_u\qelem_{k,k}=(x_k-z_1)\big(\qsch_u\qelem_{k-1,k-1}\big)+q_{k-1}\,\qsch_u\qelem_{k-2,k-2}.
\end{equation}
By the inductive hypothesis, $\qsch_u\qelem_{k-1,k-1}=\sum_v c_{k-1}(u,v)\,\qsch_v$. Applying Lemma \ref{lemma:singlevariable} with $i=k$, $j=1$ to each $\qsch_v$ and extracting the coefficient of $\qsch_w$ from \eqref{eq:pieriexpand}, we obtain
\begin{align}
[\qsch_w]\big(\qsch_u\qelem_{k,k}\big)
&=(y_{w(k)}-z_1)\,c_{k-1}(u,w) \label{eq:diag}\\
&\quad+\sum_{\substack{b>k,\ w t_{kb}\lup w\\ \ell(w)=\ell(wt_{kb})+1}} c_{k-1}(u,wt_{kb})
+\sum_{\substack{b>k,\ w t_{kb}\lup w\\ \ell(w)=\ell(wt_{kb})-2(b-k)+1}} q_{kb}\,c_{k-1}(u,wt_{kb}) \label{eq:up}\\
&\quad-\sum_{\substack{a<k,\ w t_{ak}\lup w\\ \ell(w)=\ell(wt_{ak})+1}} c_{k-1}(u,wt_{ak})
-\sum_{\substack{a<k,\ w t_{ak}\lup w\\ \ell(w)=\ell(wt_{ak})-2(k-a)+1}} q_{ak}\,c_{k-1}(u,wt_{ak}) \label{eq:down}\\
&\quad+q_{k-1}\,c_{k-2}(u,w), \label{eq:rec}
\end{align}
where we used that $\mathrm{sign}(b-k)=1$ for $b>k$ and $\mathrm{sign}(a-k)=-1$ for $a<k$, that the term $\qsch_w$ arises from $\qsch_v$ with $v=wt$ exactly when $v\lup w$, and that a summand in \eqref{eq:up}--\eqref{eq:down} is present precisely when $u\tom{k-1}wt_{kb}$, respectively $u\tom{k-1}wt_{ak}$ (otherwise $c_{k-1}=0$). In the terminology introduced before Lemma \ref{lemma:neighbors}, \eqref{eq:up} and \eqref{eq:down} are the Monk summands at $w$, with $\kappa_{kb}$, respectively $\kappa_{ak}$, equal to $1$ in the first sum and to $q_{kb}$, respectively $q_{ak}$, in the second. We must show that \eqref{eq:diag}--\eqref{eq:rec} equals $c_k(u,w)$. We distinguish three cases according to whether $u(k)=w(k)$ and, when $u(k)\neq w(k)$, whether the cycle $C$ of $u^{-1}w$ containing $k$ has its top $B$ greater than $k$ or equal to $k$.

\emph{Case 1: $u(k)=w(k)$.} By Lemma \ref{lemma:leveldichotomy}(i) we have $u\tom{k}w$ if and only if $u\tom{k-1}w$; suppose first that both hold. Then, by the same lemma,
$$\fx_{k}(u,w)=\fx_{{k-1}}(u,w)\cup\{u(k)\},\qquad q_k(u,w)=q_{k-1}(u,w),$$
so the diagonal contribution \eqref{eq:diag} gives
$$(y_{u(k)}-z_1)\,c_{k-1}(u,w)=q_{k-1}(u,w)\,(y_{u(k)}-z_1)\,\wte(y_{\fx_{{k-1}}(u,w)};z)=q_k(u,w)\,\wte(y_{\fx_{k}(u,w)};z)=c_k(u,w),$$
by Lemma \ref{lemma:elemrecursion}. The remaining contributions are the reflection sums \eqref{eq:up}, \eqref{eq:down} and the quantum term $q_{k-1}c_{k-2}(u,w)$ of \eqref{eq:rec}, which sum to $0$ by Lemma \ref{lemma:reflectioncancel}. Hence the total is $c_k(u,w)$, as required. If instead $u(k)=w(k)$ but $u\ntom{k}w$ (equivalently $u\ntom{k-1}w$), then $c_k(u,w)=0$ and the diagonal term vanishes, while the reflection sums and quantum term again cancel by Lemma \ref{lemma:reflectioncancel}; the total is $0=c_k(u,w)$.

\emph{Case 2: $u(k)\neq w(k)$ and $B>k$.} By Lemma \ref{lemma:movedcancel}(a), the diagonal term \eqref{eq:diag} and the quantum term \eqref{eq:rec} vanish, and the Monk summands \eqref{eq:up}--\eqref{eq:down} sum to $c_k(u,w)$. (When $u\tom{k}w$, Lemma \ref{lemma:pieriknotkn1} identifies the single surviving summand: it is indexed by $t_{kB}$, with $v=wt_{kB}$, $u\tom{k-1}v$, $\fx_k(u,w)=\fx_{k-1}(u,v)$ and $\kappa_{kB}\,q_{k-1}(u,v)=q_k(u,w)$, the factor $\kappa_{kB}$ being $1$ or $q_{kB}$ according as $t_{kB}$ is a classical or a quantum cover of $v$; this is the step at which the quantum drop terms of Lemma \ref{lemma:singlevariable} and the quantum weight of Definition \ref{definition:pieri} interlock.)

\emph{Case 3: $u(k)\neq w(k)$ and $B=k$.} Here $u\ntom{k}w$, since $C$ contains no element greater than $k$, so $c_k(u,w)=0$; and \eqref{eq:diag}--\eqref{eq:rec} vanishes by Lemma \ref{lemma:movedcancel}(b) (equivalently, by Lemma \ref{lemma:vanishing}).

By Lemma \ref{lemma:leveldichotomy} these cases are exhaustive. In all cases $[\qsch_w](\qsch_u\qelem_{k,k})=c_k(u,w)$, completing the induction and the proof.
\end{proof}

\section{Additional formulas}
\newcommand{\dtom}[1]{\overset{#1}{\rightrightarrows}}
\begin{definition}
	Consider a permutation $d_2[k]$ satisfying
	$$\code(d_2[k]) = \overbrace{(2,2,\ldots,2)}^{k\text{ times}}.$$
	The double Schubert polynomial $\sch_{d_2[k]}(x;y)$ is a product of two factorial elementary symmetric polynomials, specifically:
	$$\sch_{d_2[k]}(x;y) = E_k(x;y_1)E_k(x;y_2).$$
	Therefore, the Pieri formula in \cite{samuelmolev} easily transfers to being able to multiply by $\sch_{d_2[k]}(x;y)$.
	
	For quantum double Schubert polynomials, $\qsch_{d_2[k]}(x;y)$ is \emph{not} a product of factorial quantum elementary polynomials. For example:
	$$\qsch_{d_2[1]}(x;y) = \qelem_{1,2}(x;y)\qelem_{1,1}(x;y) - E^q_{2,2}(x;y)$$
	$$\qsch_{d_2[2]}(x;y) = E^q_{2,3}(x;y)E^q_{2,2}(x;y) - E^q_{3,3}(x;y)E^q_{1,2}(x;y)$$

	For abbreviation purposes, we denote this polynomial by $D_k^q(x;y)$.

	
	We define a relation $\dtom{k}$ on permutations as follows. We say that $u\dtom{k} w$ if and only if there exists a $v\in S_\infty$ such that $u\tom{k} v$ and $v\tom{k} w$ satisfying the additional condition that if $u^{-1}v$ is a product of disjoint cycles $c_1,\ldots,c_m$ and $v^{-1}w$ is a product of disjoint cycles $d_1,\ldots,d_n$, then for each $d_i$ there exists a $c_j$ such that $d_i\cap c_j = \{\max(d_i)\} = \{\max(c_j)\}$. Call the set of all such $v$ $\mathrm{pivot}_k(u,w)$.
\end{definition}

Then
\begin{theorem}[Doubled quantum Pieri formula] \label{theorem:doubledpieri}
	We have that the product $\qsch_u(x;y)D_k^q(x;z)$ is equal to
	
$$\sum_{u\dtom{k} w} \sum_{v\in\mathrm{pivot}_k(u,w)}q_k(u,v)q_k(v,w)
\prod_{a\in\fx_k(u,v)}(y_a-z_1)
\prod_{b\in\fx_k(v,w)}(y_b-z_2)\qsch_w(x;y).$$
\end{theorem}

\begin{corollary}[Iterated quantum Pieri formula] \label{corollary:iteratedpieri}
For all $u\in S_\infty$ and $k\geq 1$,
\[
\qsch_u(x;y)\qelem_{k,k}(x;z_1)\qelem_{k,k}(x;z_2)
=\sum_{u\tom{k}v\tom{k}w}
q_k(u,v)q_k(v,w)
\prod_{a\in\fx_k(u,v)}(y_a-z_1)
\prod_{b\in\fx_k(v,w)}(y_b-z_2)\qsch_w(x;y).
\]
\end{corollary}
\begin{proof}
Apply Theorem~\ref{theorem:pieri} to multiply $\qsch_u(x;y)$ by
$\qelem_{k,k}(x;z_1)$. For each resulting term indexed by $v$, apply the
same theorem with the $z$-alphabet replaced by $z_2$. The coefficient
factors multiply, and the two successive relations are exactly
$u\tom{k}v\tom{k}w$.
\end{proof}

\section{Peterson-Woodward comparison formula}

\begin{proposition}[{\cite[Proposition~2.10]{huang2015equivariant}}]
Let $W_P$ be a parabolic subgroup of $W$. Then for any $\lambda_P\in Q^\vee/Q_P^\vee$, we have that
$$c_{u,v}^{w,\lambda_P}(y;y) = c_{u,v}^{w\omega_P\omega_{P'},\lambda_B}(y;y),$$
where $\omega_P$ is the longest element in $W_P$, $\omega_{P'}$ is the longest element in $W_{P'}$, where $P'$ is the parabolic subgroup generated by $\{s_\alpha\mid \alpha\in\Delta_P,\langle \alpha, \lambda_B\rangle=0\}$, and $\lambda_B$ is the Peterson-Woodward lift of $\lambda_P$. 
\end{proposition}

\begin{corollary}
The Peterson-Woodward comparison formula holds for $c_{u,v}^w(y;z)$.
\end{corollary}
\begin{proof}
We have that
$$c_{u,v}^{w,\lambda_P}(y;z) = \sum_{v_1\cdot v_2=v}c_{u,v_2}^{w,\lambda_P}(y;y)\sch_{v_1}(y;z)$$
By definition of a parabolic subgroup, $v_2\in W_P$ as well. Thus, this is equal to
$$\sum_{v_1\cdot v_2=v}c_{u,v_2}^{w\omega_P\omega_{P'},\lambda_B}(y;y)\sch_{v_1}(y;z) = c_{u,v}^{w\omega_P\omega_{P'},\lambda_B}(y;z).$$
\end{proof}


\bibliographystyle{acm}
\bibliography{formschub}
\end{document}
